\subsection{ROOTS OF A SPLIT SEMI-SIMPLE LIE ALGEBRA}

In this number, \((\mathfrak{g},\mathfrak{h})\) denotes a split semi-simple Lie algebra.

For any \(\lambda \in \mathfrak{h}^*\) , denote by \(\mathfrak{g}^{\lambda}(\mathfrak{h})\) , or simply by \(\mathfrak{g}^{\lambda}\) , the primary subspace of \(\mathfrak{g}\) relative to \(\lambda\) (cf.~Chap. VII, §1, no. 3). Recall that \(\mathfrak{g}^{0} = \mathfrak{h}\) (Chap. VII, §2, no. 1, Prop. 4), that \(\mathfrak{g}\) is the direct sum of the \(\mathfrak{g}^{\lambda}\) (Chap. VII, §1, no. 3, Prop. 8 and 9), that \(\mathfrak{g}^{\lambda}\) is the set of \(x \in \mathfrak{g}\) such that \([h, x] = \lambda(h)x\) for all \(h \in \mathfrak{h}\) (Chap. VII, §2, no. 4, Cor. 1 of Th. 2), and that the weights of \(\mathfrak{h}\) on \(\mathfrak{g}\) are the linear forms \(\lambda\) on \(\mathfrak{h}\) such that \(\mathfrak{g}^{\lambda} \neq 0\) (Chap. VII, §1, no. 1).

DEFINITION 2. A root of \((\mathfrak{g}, \mathfrak{h})\) is a non-zero weight of h on g.

Denote by \(R(\mathfrak{g},\mathfrak{h})\) , or simply by R, the set of roots of \((\mathfrak{g},\mathfrak{h})\) . We have

\[
\mathfrak {g} = \mathfrak {h} \oplus \bigoplus_ {\alpha \in R} \mathfrak {g} ^ {\alpha}.
\]

PROPOSITION 1. Let \(\alpha, \beta\) be roots of \((\mathfrak{g}, \mathfrak{h})\) and let \(\langle \cdot, \cdot \rangle\) be a non-degenerate invariant symmetric bilinear form on \(\mathfrak{g}\) (for example the Killing form of \(\mathfrak{g}\) ).

\begin{enumerate}
\def\labelenumi{(\roman{enumi})}
\item
  If \(\alpha + \beta \neq 0\) , \(\mathfrak{g}^{\alpha}\) and \(\mathfrak{g}^{\beta}\) are orthogonal. The restriction of \(\langle \cdot, \cdot \rangle\) to \(\mathfrak{g}^{\alpha} \times \mathfrak{g}^{-\alpha}\) is non-degenerate. The restriction of \(\langle \cdot, \cdot \rangle\) to \(\mathfrak{h}\) is non-degenerate.
\item
  Let \(x \in \mathfrak{g}^{\alpha}\) , \(y \in \mathfrak{g}^{-\alpha}\) and \(h \in \mathfrak{h}\) . Then \([x, y] \in \mathfrak{h}\) and
\end{enumerate}

\[
\langle h, [ x, y ] \rangle = \alpha (h) \langle x, y \rangle .
\]

Assertion (i) is a particular case of Prop. 10 (iii) of Chap. VII, §1, no. 3. If \(x \in \mathfrak{g}^{\alpha}\) , \(y \in \mathfrak{g}^{-\alpha}\) and \(h \in \mathfrak{h}\) , we have \([x, y] \in \mathfrak{g}^{\alpha - \alpha} = \mathfrak{h}\) , and

\[
\langle h, [ x, y ] \rangle = \langle [ h, x ], y \rangle = \langle \alpha (h) x, y \rangle = \alpha (h) \langle x, y \rangle .
\]

THEOREM 1. Let \(\alpha\) be a root of \((\mathfrak{g},\mathfrak{h})\) .

\begin{enumerate}
\def\labelenumi{(\roman{enumi})}
\item
  The vector space \(\mathfrak{g}^{\alpha}\) is of dimension 1.
\item
  The vector subspace \(\mathfrak{h}_{\alpha} = [\mathfrak{g}^{\alpha},\mathfrak{g}^{-\alpha}]\) of \(\mathfrak{h}\) is of dimension 1. It contains a unique element \(H_{\alpha}\) such that \(\alpha (H_{\alpha}) = 2\) .
\item
  The vector subspace \(\mathfrak{s}_{\alpha} = \mathfrak{h}_{\alpha} + \mathfrak{g}^{\alpha} + \mathfrak{g}^{-\alpha}\) is a Lie subalgebra of \(\mathfrak{g}\) .
\item
  If \(X_{\alpha}\) is a non-zero element of \(\mathfrak{g}^{\alpha}\) , there exists a unique \(X_{-\alpha} \in \mathfrak{g}^{-\alpha}\) such that \([X_{\alpha}, X_{-\alpha}] = -H_{\alpha}\) . Let \(\varphi\) be the linear map from \(\mathfrak{sl}(2, k)\) to \(\mathfrak{g}\) that takes \(X_{+}\) to \(X_{\alpha}\) , \(X_{-}\) to \(X_{-\alpha}\) , and \(H\) to \(H_{\alpha}\) ; then \(\varphi\) is an isomorphism from the Lie algebra \(\mathfrak{sl}(2, k)\) to the Lie algebra \(\mathfrak{s}_{\alpha}\) .
\end{enumerate}

\begin{enumerate}
\def\labelenumi{\alph{enumi})}
\item
  Let \(h_{\alpha}\) be the unique element of h such that \(\alpha(h) = \langle h_{\alpha}, h \rangle\) for all \(h \in h\) . By Prop. 1, \([x, y] = \langle x, y \rangle h_{\alpha}\) for all \(x \in g^{\alpha}\) , \(y \in g^{-\alpha}\) ; on the other hand \(\langle g^{\alpha}, g^{-\alpha} \rangle \neq 0\) . Hence \(h_{\alpha} = [g^{\alpha}, g^{-\alpha}] = kh_{\alpha}\) .
\item
  Choose \(x \in \mathfrak{g}^{\alpha}\) , \(y \in \mathfrak{g}^{-\alpha}\) such that \(\langle x, y \rangle = 1\) , so \([x, y] = h_{\alpha}\) . Recall that \([h_{\alpha}, x] = \alpha(h_{\alpha})x\) , \([h_{\alpha}, y] = -\alpha(h_{\alpha})y\) . If \(\alpha(h_{\alpha}) = 0\) , it follows that \(kx + ky + kh_{\alpha}\) is a nilpotent subalgebra \(t\) of \(\mathfrak{g}\) ; since \(h_{\alpha} \in [t, t]\) , \(\operatorname{ad}_{\mathfrak{g}} h_{\alpha}\) is nilpotent (Chap. I, §5, no. 3, Th. 1), which is absurd since \(\operatorname{ad}_{\mathfrak{g}} h_{\alpha}\) is non-zero semi-simple. So \(\alpha(h_{\alpha}) \neq 0\) . Hence there exists a unique \(H_{\alpha} \in \mathfrak{h}_{\alpha}\) such that \(\alpha(H_{\alpha}) = 2\) , which proves (ii).
\item
  Choose a non-zero element \(X_{\alpha}\) of \(\mathfrak{g}^{\alpha}\) . There exists \(X_{-\alpha} \in \mathfrak{g}^{-\alpha}\) such that \([X_{\alpha}, X_{-\alpha}] = -H_{\alpha}\) (since \([X_{\alpha}, \mathfrak{g}^{-\alpha}] = \mathfrak{h}_{\alpha}\) by \(b\) ). Then
\end{enumerate}

\[
\begin{array}{c} [ H _ {\alpha}, X _ {\alpha} ] = \alpha (H _ {\alpha}) X _ {\alpha} = 2 X _ {\alpha}, [ H _ {\alpha}, X _ {- \alpha} ] = - \alpha (H _ {\alpha}) X _ {- \alpha} = - 2 X _ {- \alpha}, \\ \relax [ X _ {\alpha}, X _ {- \alpha} ] = - H _ {\alpha}; \end{array}
\]

hence \(kX_{\alpha} + kX_{-\alpha} + kH_{\alpha}\) is a subalgebra of \(\mathfrak{g}\) and the linear map \(\varphi\) from \(\mathfrak{sl}(2,k)\) to \(kX_{\alpha} + kX_{-\alpha} + kH_{\alpha}\) such that \(\varphi(X_{+}) = X_{\alpha}, \varphi(X_{-}) = X_{-\alpha}, \varphi(H) = H_{\alpha}\) is an isomorphism of Lie algebras.

\begin{enumerate}
\def\labelenumi{\alph{enumi})}
\setcounter{enumi}{3}
\tightlist
\item
  Assume that \(\dim g^{\alpha} > 1\) . Let y be a non-zero element of \(g^{-\alpha}\) . There exists a non-zero element \(X_{\alpha}\) of \(g_{\alpha}\) such that \(\langle y, X_{\alpha} \rangle = 0\) . Choose \(X_{-\alpha}\) as in c), and consider the representation \(\rho : u \mapsto \operatorname{ad}_{\mathfrak{g}} \varphi(u)\) from \(\mathfrak{sl}(2, k)\) to g. We have
\end{enumerate}

\[
\begin{array}{c} \rho (H) y = [ \varphi (H), y ] = [ H _ {\alpha}, y ] = - 2 y \\ \rho (X _ {+}) y = [ \varphi (X _ {+}), y ] = [ X _ {\alpha}, y ] = \langle X _ {\alpha}, y \rangle h _ {\alpha} = 0. \end{array}
\]

Thus, \(y\) is primitive for \(\rho\) , of weight \(-2\) , which contradicts Prop. 2 of §1, no. 2. This proves (i).

\begin{enumerate}
\def\labelenumi{\alph{enumi})}
\setcounter{enumi}{4}
\tightlist
\item
  Assertion (iii) is now a consequence of c). On the other hand, if \(X_{\alpha}\) is a non-zero element of \(g^{\alpha}\) , the element \(X_{-\alpha}\) constructed in c) is the unique element of \(g^{-\alpha}\) such that \([X_{\alpha}, X_{-\alpha}] = -H_{\alpha}\) since \(\dim g^{-\alpha} = 1\) . The last assertion of (iv) is a consequence of c). Q.E.D.
\end{enumerate}

The notations \(h_{\alpha}\) , \(H_{\alpha}\) , \(s_{\alpha}\) will be retained in what follows. (To define \(h_{\alpha}\) , we take \(\langle\cdot,\cdot\rangle\) equal to the Killing form.) If \(X_{\alpha}\) is a non-zero element of \(g^{\alpha}\) , the isomorphism \(\varphi\) of Th. 1 and the representation \(u \mapsto \operatorname{ad}_{\mathfrak{g}} \varphi(u)\) of \(\mathfrak{sl}(2,k)\) on g will be said to be associated to \(X_{\alpha}\) .

COROLLARY. Let \(\Phi\) be the Killing form of \(\mathfrak{g}\) . For all \(a, b \in \mathfrak{h}\) ,

\[
\Phi (a, b) = \sum_ {\gamma \in \mathrm{R}} \gamma (a) \gamma (b).
\]

Indeed, ad a.ad b leaves each \(g^{\gamma}\) stable, and its restriction to \(g^{\gamma}\) is the homothety with ratio \(\gamma(a)\gamma(b)\) ; if \(\gamma\neq0\) , \(\dim g^{\gamma}=1\) .

PROPOSITION 2. Let \(\alpha, \beta \in \mathbb{R}\) .

\begin{enumerate}
\def\labelenumi{(\roman{enumi})}
\item
  \(\beta(H_{\alpha}) \in \mathbf{Z}\) .
\item
  If \(\Phi\) denotes the Killing form of \(\mathfrak{g}\) , \(\Phi(H_{\alpha}, H_{\beta}) \in \mathbf{Z}\) .
\end{enumerate}

Let \(X_{\alpha}\) be a non-zero element of \(\mathfrak{g}^{\alpha}\) , and let \(\rho\) be the representation of \(\mathfrak{sl}(2,k)\) on \(\mathfrak{g}\) associated to \(X_{\alpha}\) . The eigenvalues of \(\rho(\mathrm{H})\) are 0 and the \(\beta(H_{\alpha})\)

for \(\beta \in \mathbb{R}\) . Hence (i) follows from §1, no. 2, Cor. of Prop. 2. Assertion (ii) follows from (i) and the Cor. of Th. 1. Q.E.D.

Let \(\alpha \in \mathrm{R}\) , \(X_{\alpha}\) a non-zero element of \(\mathfrak{g}^{\alpha}\) , \(X_{-\alpha}\) the element of \(\mathfrak{g}^{-\alpha}\) such that \([X_{\alpha}, X_{-\alpha}] = -H_{\alpha}\) , and \(\rho\) the representation of \(\mathfrak{sl}(2, k)\) on \(\mathfrak{g}\) associated to \(X_{\alpha}\) . Let \(\pi\) be the representation of \(\mathbf{SL}(2, k)\) on \(\mathfrak{g}\) compatible with \(\rho\) ( \(\S 1\) , no. 4, Th. 2). Since \(\operatorname{ad} X_{\alpha}\) is nilpotent (Chap. VII, \(\S 1\) , no. 3, Prop. 10 (iv)), \(\pi(e^{X_+}) = e^{\operatorname{ad} X_{\alpha}}\) is an elementary automorphism of \(\mathfrak{g}\) . Similarly, \(\pi(e^{X_-}) = e^{\operatorname{ad} X_{-\alpha}}\) is an elementary automorphism of \(\mathfrak{g}\) . Hence \(\pi(\mathbf{SL}(2, k)) \subset \operatorname{Aut}_e(\mathfrak{g})\) . We make use of the notation \(\theta(t)\) of \(\S 1\) , no. 5. For \(t \in k^*\) , put

\[
\theta_ {\alpha} (t) = \pi (\theta (t)) = e ^ {\mathrm{ad} t X _ {\alpha}} e ^ {\mathrm{ad} t ^ {- 1} X _ {- \alpha}} e ^ {\mathrm{ad} t X _ {\alpha}}.\tag{1}
\]

Lemma 1. (i) For all \(h \in \mathfrak{h}\) , \(\theta_{\alpha}(t).h = h - \alpha(h)H_{\alpha}\) .

\begin{enumerate}
\def\labelenumi{(\roman{enumi})}
\setcounter{enumi}{1}
\item
  For all \(\beta \in \mathrm{R}\) , \(\theta_{\alpha}(t)(\mathfrak{g}^{\beta}) = \mathfrak{g}^{\beta - \beta (H_{\alpha})\alpha}\) .
\item
  If \(\alpha, \beta \in \mathrm{R}\) , \(\beta - \beta(H_{\alpha})\alpha \in \mathrm{R}\) .
\end{enumerate}

Let \(h \in \mathfrak{h}\) . If \(\alpha(h) = 0\) , \([X_{\alpha}, h] = [X_{-\alpha}, h] = 0\) , so \(\theta_{\alpha}(t).h = h\) . On the other hand, the formulas (5) of §1, no. 5 show that \(\theta_{\alpha}(t).H_{\alpha} = -H_{\alpha}\) . This proves assertion (i). It follows that \(\theta_{\alpha}(t)^2 | \mathfrak{h} = \mathrm{Id}\) . If \(x \in \mathfrak{g}^{\beta}\) and \(h \in \mathfrak{h}\) ,

\[
\begin{array}{r l} [ h, \theta_ {\alpha} (t) x ] & = \theta_ {\alpha} (t). [ \theta_ {\alpha} (t) h, x ] - \beta (\theta_ {\alpha} (t) h). \theta_ {\alpha} (t) x \\ & = (\beta (h) - \alpha (h) \beta (H _ {\alpha})). \theta_ {\alpha} (t) x \\ & = (\beta - \beta (H _ {\alpha}) \alpha) (h). \theta_ {\alpha} (t) x \end{array}
\]

so \(\theta_{\alpha}(t)x\in \mathfrak{g}^{\beta -\beta (H_{\alpha})\alpha}\) . This proves (ii). Assertion (iii) follows from (ii).

THEOREM 2. (i) The set \(\mathrm{R} = \mathrm{R}(\mathfrak{g},\mathfrak{h})\) is a reduced root system in \(\mathfrak{h}^*\) .

\begin{enumerate}
\def\labelenumi{(\roman{enumi})}
\setcounter{enumi}{1}
\tightlist
\item
  Let \(\alpha \in \mathbb{R}\) . The map \(s_{\alpha, H_\alpha}: \lambda \mapsto \lambda - \lambda(H_\alpha)\alpha\) from \(\mathfrak{h}^*\) to \(\mathfrak{h}^*\) is the unique reflection \(s\) of \(\mathfrak{h}^*\) such that \(s(\alpha) = -\alpha\) and \(s(\mathrm{R}) = \mathrm{R}\) . For all \(t \in k^*\) , \(s\) is the transpose of \(\theta_\alpha(t)|\mathfrak{h}\) .
\end{enumerate}

First, R generates \(h^{*}\) , for if \(h \in h\) is such that \(\alpha(h) = 0\) for all \(\alpha \in R\) , then ad h = 0 and hence h = 0 since the centre of g is zero. By definition, \(0 \notin R\) . Let \(\alpha \in R\) . Since \(\alpha(H_{\alpha}) = 2\) , \(s = s_{\alpha,H_{\alpha}}\) is a reflection such that \(s(\alpha) = -\alpha\) . Then \(s(R) = R\) by Lemma 1 (iii), and \(\beta(H_{\alpha}) \in \mathbf{Z}\) for all \(\beta \in R\) (Prop. 2 (i)). This shows that R is a root system in \(h^{*}\) . For all \(h \in h\) and \(\lambda \in h^{*}\) ,

\[
\langle s (\lambda), h \rangle = \langle \lambda - \lambda (H _ {\alpha}) \alpha , h \rangle = \langle \lambda , h - \alpha (h) H _ {\alpha} \rangle = \langle \lambda , \theta_ {\alpha} (t) h \rangle
\]

so \(s\) is the transpose of \(\theta_{\alpha}(t)|\mathfrak{h}\) . Finally, we show that the root system R is reduced. Let \(\alpha \in \mathrm{R}\) and \(y \in \mathfrak{g}^{2\alpha}\) . Since \(3\alpha \notin \mathrm{R}\) (Chap. VI, §1, no. 3, Prop. 8), \([X_{\alpha}, y] = 0\) ; on the other hand, \([X_{-\alpha}, y] \in \mathfrak{g}^{-\alpha + 2\alpha} = \mathfrak{g}^{\alpha} = kX_{\alpha}\) , so \([X_{\alpha}, [X_{-\alpha}, y]] = 0\) ; thus

\[
4 y = 2 \alpha (H _ {\alpha}) y = [ H _ {\alpha}, y ] = - [ [ X _ {\alpha}, X _ {- \alpha} ], y ] = 0
\]

so \(y = 0\) and \(\mathfrak{g}^{2\alpha} = 0\) . In other words, \(2\alpha\) is not a root.

Q.E.D.

Identify \(\mathfrak{h}\) canonically with \(\mathfrak{h}^{**}\) . With the notations of Chap. VI, §1, no. 1, we then have, by Th. 2 (ii),

\[
H _ {\alpha} = \alpha^ {\vee} \quad \text { for   all } \alpha \in \mathrm{R}.\tag{2}
\]

The \(H_{\alpha}\) thus form the root system \(\mathbb{R}^{\vee}\) in \(\mathfrak{h}\) inverse to \(\mathbb{R}\) .

We shall call \(\mathrm{R}(\mathfrak{g},\mathfrak{h})\) the root system of \((\mathfrak{g},\mathfrak{h})\) . The reflections \(s_{\alpha,H_{\alpha}}\) will be denoted simply by \(s_{\alpha}\) . The Weyl group, group of weights, Coxeter number \ldots{} of \(\mathrm{R}(\mathfrak{g},\mathfrak{h})\) are called the Weyl group, group of weights, Coxeter number \ldots{} of \((\mathfrak{g},\mathfrak{h})\) . As in Chap. VI, §1, no. 1, we consider the Weyl group as operating not only on \(h^{*}\) , but also on h by transport of structure, so that \(s_{\alpha} = \theta_{\alpha}(t)|\mathfrak{h}\) . Since the \(\theta_{\alpha}(t)\) are elementary automorphisms of g, we have:

COROLLARY. Every element of the Weyl group of \((\mathfrak{g},\mathfrak{h})\) , operating on h, is the restriction to h of an elementary automorphism of g.

For a converse of this result, see §5, no. 2, Prop. 4.

Remark 1. If \(\mathfrak{h}_{\mathbf{Q}}\) (resp. \(\mathfrak{h}_{\mathbf{Q}}^{*}\) ) denotes the \(\mathbf{Q}\) -vector subspace of \(\mathfrak{h}\) (resp. \(\mathfrak{h}^{*}\) ) generated by the \(H_{\alpha}\) (resp. the \(\alpha\) ), where \(\alpha \in \mathbb{R}\) , then \(\mathfrak{h}\) (resp. \(\mathfrak{h}^{*}\) ) can be identified canonically with \(\mathfrak{h}_{\mathbf{Q}} \otimes_{\mathbf{Q}} k\) (resp. with \(\mathfrak{h}_{\mathbf{Q}}^{*} \otimes_{\mathbf{Q}} k\) ) and \(\mathfrak{h}_{\mathbf{Q}}^{*}\) can be identified with the dual of \(\mathfrak{h}_{\mathbf{Q}}\) (Chap. VI, §1, no. 1, Prop. 1). We call \(\mathfrak{h}_{\mathbf{Q}}\) and \(\mathfrak{h}_{\mathbf{Q}}^{*}\) the canonical \(\mathbf{Q}\) -structures on \(\mathfrak{h}\) and \(\mathfrak{h}^{*}\) (Algebra, Chap. II, §8, no. 1, Def. 1). When we mention \(\mathbf{Q}\) -rationality for a vector subspace of \(\mathfrak{h}\) , for a linear form on \(\mathfrak{h}\) , etc., we shall mean these structures, unless we indicate otherwise. When we mention Weyl chambers, or facets, of R(g, h), we shall work in \(\mathfrak{h}_{\mathbf{Q}} \otimes_{\mathbf{Q}} \mathbf{R}\) or \(\mathfrak{h}_{\mathbf{Q}}^{*} \otimes_{\mathbf{Q}} \mathbf{R}\) , that we shall denote by \(\mathfrak{h}_{\mathbf{R}}\) and \(\mathfrak{h}_{\mathbf{R}}^{*}\) .

Remark 2. The root system \(\mathbb{R}^{\vee}\) in \(\mathfrak{h}\) defines a non-degenerate symmetric bilinear form \(\beta\) on \(\mathfrak{h}\) (Chap. VI, §1, no. 1, Prop. 3), namely the form \((a,b)\mapsto \sum_{\alpha \in \mathbb{R}}\langle \alpha ,a\rangle \langle \alpha ,b\rangle\) . By the Cor. to Th. 1, this form is just the restriction of the Killing form to \(\mathfrak{h}\) . The extension of \(\beta |\mathfrak{h}_{\mathbf{Q}}\times \mathfrak{h}_{\mathbf{Q}}\) to \(\mathfrak{h}_{\mathbf{Q}}\otimes_{\mathbf{Q}}\mathbf{R}\) is positive nondegenerate (Chap. VI, §1, no. 1, Prop. 3). On the other hand, we see that the inverse form on \(\mathfrak{h}^*\) of the restriction to \(\mathfrak{h}\) of the Killing form on \(\mathfrak{g}\) is the canonical bilinear form \(\Phi_{\mathrm{R}}\) of R (Chap. VI, §1, no. 12).

Let \((\mathfrak{g}_{1},\mathfrak{h}_{1})\) , \((\mathfrak{g}_{2},\mathfrak{h}_{2})\) be split semi-simple Lie algebras, \(\varphi\) an isomorphism from \(\mathfrak{g}_{1}\) to \(\mathfrak{g}_{2}\) such that \(\varphi(\mathfrak{h}_{1})=\mathfrak{h}_{2}\) . By transport of structure, the transpose of the map \(\varphi|\mathfrak{h}_{1}\) takes \(R(\mathfrak{g}_{2},\mathfrak{h}_{2})\) to \(R(\mathfrak{g}_{1},\mathfrak{h}_{1})\) .

PROPOSITION 3. Let g be a semi-simple Lie algebra, \(h_{1}\) and \(h_{2}\) splitting Cartan subalgebras of g. There exists an isomorphism from \(h_{1}^{*}\) to \(h_{2}^{*}\) that takes \(\mathrm{R}(\mathfrak{g},\mathfrak{h}_{1})\) to \(\mathrm{R}(\mathfrak{g},\mathfrak{h}_{2})\) .

(For more precise results, see §3, no. 3, Cor. of Prop. 10, and §5, no. 3, Prop. 5).

Let \(k'\) be an algebraic closure of \(k\) , \(\mathfrak{g}' = \mathfrak{g} \otimes_k k'\) , \(\mathfrak{h}_i' = \mathfrak{h}_i \otimes_k k'\) . Then \(\mathrm{R}(\mathfrak{g}', \mathfrak{h}_i')\) is the image of \(\mathrm{R}(\mathfrak{g}, \mathfrak{h}_i)\) under the map \(\lambda \mapsto \lambda \otimes 1\) from \(\mathfrak{h}_i^*\) to \(\mathfrak{h}_i^* \otimes_k k' = \mathfrak{h}_i'^*\) . By Chap. VII, §3, no. 2, Th. 1, there exists an automorphism of \(\mathfrak{g}'\) taking \(\mathfrak{h}_1'\) to \(\mathfrak{h}_2'\) , hence an isomorphism \(\varphi\) from \(\mathfrak{h}_1'^*\) to \(\mathfrak{h}_2'^*\) that takes \(\mathrm{R}(\mathfrak{g}', \mathfrak{h}_1')\) to \(\mathrm{R}(\mathfrak{g}', \mathfrak{h}_2')\) . Then \(\varphi | \mathfrak{h}_1^*\) takes \(\mathrm{R}(\mathfrak{g}, \mathfrak{h}_1)\) to \(\mathrm{R}(\mathfrak{g}, \mathfrak{h}_2)\) , and hence \(\mathfrak{h}_1^*\) to \(\mathfrak{h}_2^*\) . Q.E.D.

In view of Prop. 3, the root system of \((\mathfrak{g},\mathfrak{h})\) depends, up to isomorphism, only on g and not on h. In the same way, the Weyl group, group of weights \ldots{} of \((\mathfrak{g},\mathfrak{h})\) are simply called, by abuse of language, the Weyl group, group of weights \ldots{} of g (cf.~also §5, no. 3, Remark 2). If the Dynkin graph of g is of type \(A_{l}\) , or \(B_{l}\) , \ldots{} (cf.~Chap. VI, §4, no. 2, Th. 3), we say that g is of type \(A_{l}\) , or \(B_{l}\) , \ldots.

Recall that, if \(\alpha\) and \(\beta\) are linearly independent roots, the set of \(j \in \mathbf{Z}\) such that \(\beta + j\alpha \in \mathrm{R}\) is an interval \([-q, p]\) of \(\mathbf{Z}\) containing 0, with \(p - q = -\langle \beta, \alpha^{\vee} \rangle = -\beta(H_{\alpha})\) (Chap. VI, §1, no. 3, Prop. 9).

PROPOSITION 4. Let \(\alpha\) and \(\beta\) be linearly independent roots. Let \(p\) (resp. \(q\) ) be the largest integer \(j\) such that \(\beta + j\alpha\) (resp. \(\beta - j\alpha\) ) is a root.

\begin{enumerate}
\def\labelenumi{(\roman{enumi})}
\item
  The vector subspace \(\sum_{-q\leq j\leq p}\mathfrak{g}^{\beta +j\alpha}\) of \(\mathfrak{g}\) is a simple \(\mathfrak{s}_{\alpha}\) -module of dimension \(p + q + 1\) .
\item
  If \(\alpha + \beta\) is a root, then \([\mathfrak{g}^{\alpha}, \mathfrak{g}^{\beta}] = \mathfrak{g}^{\alpha + \beta}\) .
\end{enumerate}

Let \(X_{\alpha}\) (resp. x) be a non-zero element of \(\mathfrak{g}^{\alpha}\) (resp. \(\mathfrak{g}^{\beta+p\alpha}\) ). Then

\[
\begin{array}{c} {[ X _ {\alpha}, x ] \in \mathfrak {g} ^ {\beta + (p + 1) \alpha} = 0} \\ {[ H _ {\alpha}, x ] = (\beta (H _ {\alpha}) + p \alpha (H _ {\alpha})) x = (- p + q + 2 p) x = (p + q) x.} \end{array}
\]

Thus, \(x\) is primitive of weight \(p + q\) for the representation of \(\mathfrak{sl}(2,k)\) on \(\mathfrak{g}\) associated to \(X_{\alpha}\) ; but the \(\mathfrak{sl}(2,k)\) -module \(\sum_{-q\leq j\leq p}\mathfrak{g}^{\beta +j\alpha}\) is of dimension \(p + q + 1\) ; hence it is simple (§1, no. 2, Prop. 2). If \(\alpha +\beta \in \mathrm{R}\) , then \(p\geq 1\) , so the elements of \(\mathfrak{g}^{\beta}\) are not primitive, and hence \([X_{\alpha},\mathfrak{g}^{\beta}]\neq 0\) . Since \([\mathfrak{g}^{\alpha},\mathfrak{g}^{\beta}]\subset \mathfrak{g}^{\alpha +\beta}\) , we see finally that \([\mathfrak{g}^{\alpha},\mathfrak{g}^{\beta}] = \mathfrak{g}^{\alpha +\beta}\) .

Remark 3. Recall that, by Chap. VI, §1, no. 3, Cor. of Prop. 9, the integer \(p + q + 1\) can only take the values 1, 2, 3, 4.

Remark 4. Let \((\mathfrak{g},\mathfrak{h})\) be a split reductive Lie algebra, \(\mathfrak{c}\) the centre of \(\mathfrak{g},\mathfrak{g}' = \mathcal{D}\mathfrak{g}\) , \(\mathfrak{h}' = \mathfrak{h} \cap \mathfrak{g}'\) . Then \(\mathfrak{h} = \mathfrak{c} \times \mathfrak{h}'\) , and we identify \(\mathfrak{h}'^*\) with a vector subspace of \(\mathfrak{h}^*\) . For any \(\lambda \in \mathfrak{h}^*\) such that \(\lambda \neq 0\) , the primary subspace \(\mathfrak{g}^\lambda\) relative to \(\lambda\) is equal to \(\mathfrak{g}'^{\lambda | \mathfrak{h}'}\) . A non-zero weight of \(\mathfrak{h}\) on \(\mathfrak{g}\) is called a root of \((\mathfrak{g},\mathfrak{h})\) ; every root vanishes on \(\mathfrak{c}\) . Denote by \(\mathrm{R}(\mathfrak{g},\mathfrak{h})\) the set of roots of \((\mathfrak{g},\mathfrak{h})\) ; it can be identified canonically with \(\mathrm{R}(\mathfrak{g}',\mathfrak{h}')\) . Let \(\alpha \in \mathrm{R}(\mathfrak{g},\mathfrak{h})\) . We define \(h_\alpha\) , \(H_\alpha\) , \(\mathfrak{s}_\alpha\) , the isomorphisms \(\mathfrak{sl}(2,k) \to \mathfrak{s}_\alpha\) , and the representations of \(\mathfrak{sl}(2,k)\) on \(\mathfrak{g}\) associated to \(\alpha\) , as in the semi-simple case.

\subsection{INVARIANT BILINEAR FORMS}

PROPOSITION 5. Let \((\mathfrak{g},\mathfrak{h})\) be a split semi-simple Lie algebra, \(\Phi\) an invariant symmetric bilinear form on g, and W the Weyl group of \((\mathfrak{g},\mathfrak{h})\) . Then the restriction \(\Phi'\) of \(\Phi\) to h is invariant under W. Moreover, if \(\Phi\) is nondegenerate, so is \(\Phi'\) .

Let \(\alpha \in \mathbb{R}\) , let \(X_{\alpha}\) be a non-zero element of \(\mathfrak{g}^{\alpha}\) , \(\rho\) the associated representation of \(\mathfrak{sl}(2,k)\) on \(\mathfrak{g}\) , and \(\pi\) the representation of \(\mathbf{SL}(2,k)\) on \(\mathfrak{g}\) compatible with \(\rho\) . Then \(\varPhi\) is invariant under \(\rho\) , and hence under \(\pi\) (§1, no. 4). In particular, \(\varPhi'\) is invariant under \(\theta_{\alpha}(t)|\mathfrak{h}\) (no. 2), and hence under W. The last assertion follows from Prop. 1 (i).

PROPOSITION 6. Let \((\mathfrak{g},\mathfrak{h})\) be a split semi-simple Lie algebra, \(\varPhi\) a nondegenerate invariant symmetric bilinear form on \(\mathfrak{g}\) . For all \(\alpha \in \mathbb{R}\) , let \(X_{\alpha}\) be a non-zero element of \(\mathfrak{g}^{\alpha}\) . Let \((H_i)_{i\in I}\) be a basis of \(\mathfrak{h}\) , and \((H_i')_{i\in I}\) the basis of \(\mathfrak{h}\) such that \(\varPhi(H_i,H_j') = \delta_{ij}\) . The Casimir element associated to \(\varPhi\) in the enveloping algebra of \(\mathfrak{g}\) (Chap. I, §3, no. 7) is then

\[
\sum_ {\alpha \in \mathrm{R}} \frac {1}{\varPhi (X _ {\alpha} , X _ {- \alpha})} X _ {\alpha} X _ {- \alpha} + \sum_ {i \in I} H _ {i} H _ {i} ^ {\prime}.
\]

Indeed, by Prop. 1, \(\Phi(H_i, X_\alpha) = \Phi(H_i', X_\alpha) = 0\) for all \(i \in I\) , \(\alpha \in \mathrm{R}\) , and \(\Phi\left(\frac{1}{\Phi(X_\alpha, X_{-\alpha})} X_\alpha, X_{-\beta}\right) = \delta_{\alpha\beta}\) for all \(\alpha, \beta \in \mathrm{R}\) .

\subsection{\texorpdfstring{THE COEFFICIENTS \(\mathbf{N}_{\alpha \beta}\)}{THE COEFFICIENTS \textbackslash mathbf\{N\}\_\{\textbackslash alpha \textbackslash beta\}}}

In this number, we again denote by \((\mathfrak{g},\mathfrak{h})\) a split semi-simple Lie algebra.

Lemma 2. There exists a family \((X_{\alpha})_{\alpha \in \mathbb{R}}\) such that, for all \(\alpha \in \mathbb{R}\) ,

\[
X _ {\alpha} \in \mathfrak {g} ^ {\alpha} \text { and } [ X _ {\alpha}, X _ {- \alpha} ] = - H _ {\alpha}.
\]

Let \(R_{1}\) be a subset of R such that \(R = R_{1} \cup (-R_{1})\) and \(R_{1} \cap (-R_{1}) = \varnothing\) . For \(\alpha \in R_{1}\) , choose an arbitrary non-zero element \(X_{\alpha}\) of \(g^{\alpha}\) . There exists a unique \(X_{-\alpha} \in g^{-\alpha}\) such that \([X_{\alpha}, X_{-\alpha}] = -H_{\alpha}\) (Th. 1 (iv)). Then

\[
\left[ X _ {- \alpha}, X _ {\alpha} \right] = H _ {\alpha} = - H _ {- \alpha}.\tag{Q.E.D.}
\]

If \((X_{\alpha})_{\alpha \in \mathbb{R}}\) is one family satisfying the conditions of Lemma 2, the most general family satisfying these conditions is \((t_{\alpha}X_{\alpha})_{\alpha \in \mathbb{R}}\) where \(t_{\alpha}\in k^{*}\) and \(t_{\alpha}t_{-\alpha} = 1\) for all \(\alpha \in \mathbb{R}\) .

In the remainder of this number, we denote by \((X_{\alpha})_{\alpha\in\mathbb{R}}\) a family satisfying the conditions of Lemma 2. We denote by \(\langle\cdot,\cdot\rangle\) a non-degenerate invariant symmetric bilinear form on g.

Every \(x \in \mathfrak{g}\) can be written uniquely in the form

\[
x = h + \sum_ {\alpha \in \mathrm{R}} \mu_ {\alpha} X _ {\alpha} \quad (h \in \mathfrak {h}, \mu_ {\alpha} \in k).
\]

The bracket of two such elements can be calculated by means of the following formulas:

\[
[ h, X _ {\alpha} ] = \alpha (h) X _ {\alpha}
\]

\[
[ X _ {\alpha}, X _ {\beta} ] = \left\{ \begin{array}{l l} 0 & \text {if} \alpha + \beta \notin \mathrm{R} \cup \{0 \} \\ - H _ {\alpha} & \text {if} \alpha + \beta = 0 \\ \mathrm{N} _ {\alpha \beta} X _ {\alpha + \beta} & \text {if} \alpha + \beta \in \mathrm{R} \end{array} \right.
\]

the \(N_{\alpha\beta}\) being non-zero elements of k.

Lemma 3. For all \(\alpha \in \mathbb{R}\) ,

\[
\langle X _ {\alpha}, X _ {- \alpha} \rangle = - \frac {1}{2} \langle H _ {\alpha}, H _ {\alpha} \rangle .
\]

Indeed,

\[
\begin{array}{c} 2 \langle X _ {\alpha}, X _ {- \alpha} \rangle = \langle \alpha (H _ {\alpha}) X _ {\alpha}, X _ {- \alpha} \rangle = \langle [ H _ {\alpha}, X _ {\alpha} ], X _ {- \alpha} \rangle \\ = \langle H _ {\alpha}, [ X _ {\alpha}, X _ {- \alpha} ] \rangle = - \langle H _ {\alpha}, H _ {\alpha} \rangle . \end{array}
\]

Lemma 4. Let \(\alpha, \beta \in \mathrm{R}\) be such that \(\alpha + \beta \in \mathrm{R}\) . Let \(p\) (resp. \(q\) ) be the largest integer \(j\) such that \(\beta + j\alpha \in \mathrm{R}\) (resp. \(\beta - j\alpha \in \mathrm{R}\) ). Then,

\[
\mathrm{N} _ {\alpha , \beta} \mathrm{N} _ {- \alpha , \alpha + \beta} = - p (q + 1)\tag{3}
\]

\[
\mathrm{N} _ {- \alpha , \alpha + \beta} \langle H _ {\beta}, H _ {\beta} \rangle = - \mathrm{N} _ {- \alpha , - \beta} \langle H _ {\alpha + \beta}, H _ {\alpha + \beta} \rangle\tag{4}
\]

\[
\mathrm{N} _ {\alpha , \beta} \mathrm{N} _ {- \alpha , - \beta} = (q + 1) ^ {2}.\tag{5}
\]

Let \(\rho\) be the representation of \(\mathfrak{sl}(2,k)\) on \(\mathfrak{g}\) defined by \(X_{\alpha}\) . The element \(e = X_{\beta + p\alpha}\) is primitive of weight \(p + q\) (Prop. 4 (i)). Put

\[
e _ {n} = \frac {(- 1) ^ {n}}{n !} \rho (X _ {-}) ^ {n} e \quad \text { for } n \geq 0.
\]

By Prop. 1 of §1,

\[
(\mathrm{ad} X _ {\alpha}) e _ {p} = (q + 1) e _ {p - 1}
\]

\[
(\operatorname{ad} X _ {- \alpha}) (\operatorname{ad} X _ {\alpha}) e _ {p} = - p (q + 1) e _ {p}.
\]

This proves (3) since \(e_p\) is a non-zero element of \(\mathfrak{g}^{\beta}\) .

The form \(\langle \cdot, \cdot \rangle\) being invariant, we have

\[
\langle [ X _ {- \alpha}, X _ {\alpha + \beta} ], X _ {- \beta} \rangle = - \langle X _ {\alpha + \beta}, [ X _ {- \alpha}, X _ {- \beta} ] \rangle
\]

SO

\[
\mathrm{N} _ {- \alpha , \alpha + \beta} \langle X _ {\beta}, X _ {- \beta} \rangle = - \mathrm{N} _ {- \alpha , - \beta} \langle X _ {\alpha + \beta}, X _ {- \alpha - \beta} \rangle
\]

which, in view of Lemma 3, proves (4).

The restriction of \(\langle\cdot,\cdot\rangle\) to h is non-degenerate and invariant under the Weyl group (Prop. 5). Identify h and \(h^{*}\) by means of this restriction. If \(\gamma\in R\) , \(H_{\gamma}\) is identified with \(2\gamma/\langle\gamma,\gamma\rangle\) (Chap. VI, §1, no. 1, Lemma 2); hence, for all \(\gamma,\delta\in R\) ,

\[
\frac {\langle \gamma , \gamma \rangle}{\langle \delta , \delta \rangle} = \frac {\langle H _ {\delta} , H _ {\delta} \rangle}{\langle H _ {\gamma} , H _ {\gamma} \rangle}.\tag{6}
\]

Now, by Chap. VI, §1, no. 3, Prop. 10,

\[
\frac {\langle \alpha + \beta , \alpha + \beta \rangle}{\langle \beta , \beta \rangle} = \frac {q + 1}{p}\tag{7}
\]

so, by (3), (4), (6), (7),

\[
\begin{array}{r} \mathrm{N} _ {\alpha , \beta} \mathrm{N} _ {- \alpha , - \beta} = - \mathrm{N} _ {\alpha , \beta} \mathrm{N} _ {- \alpha , \alpha + \beta} \frac {\langle H _ {\beta} , H _ {\beta} \rangle}{\langle H _ {\alpha + \beta} , H _ {\alpha + \beta} \rangle} \\ = - \mathrm{N} _ {\alpha , \beta} \mathrm{N} _ {- \alpha , \alpha + \beta} \frac {q + 1}{p} = (q + 1) ^ {2}. \end{array}
\]

DEFINITION 3. A Chevalley system for \((\mathfrak{g},\mathfrak{h})\) is a family \((X_{\alpha})_{\alpha \in \mathbb{R}}\) such that

\begin{enumerate}
\def\labelenumi{(\roman{enumi})}
\item
  \(X_{\alpha} \in g^{\alpha}\) for all \(\alpha \in R;\)
\item
  \([X_{\alpha}, X_{-\alpha}] = -H_{\alpha}\) for all \(\alpha \in \mathrm{R}\) ;
\item
  the linear map from \(\mathfrak{g}\) to \(\mathfrak{g}\) which is equal to \(-1\) on \(\mathfrak{h}\) and which takes \(X_{\alpha}\) to \(X_{-\alpha}\) for all \(\alpha \in \mathrm{R}\) is an automorphism of \(\mathfrak{g}\) .
\end{enumerate}

The extension of this definition to the case where \((\mathfrak{g},\mathfrak{h})\) is split reductive is immediate.

We shall show (§4, no. 4, Cor. of Prop. 5) that Chevalley systems for \((\mathfrak{g},\mathfrak{h})\) exist.

PROPOSITION 7. Let \((X_{\alpha})_{\alpha \in \mathbb{R}}\) be a Chevalley system for \((\mathfrak{g},\mathfrak{h})\) . We retain the notation of Lemma 4. Then, \(\mathrm{N}_{-\alpha, - \beta} = \mathrm{N}_{\alpha ,\beta}\) and \(\mathrm{N}_{\alpha ,\beta} = \pm (q + 1)\) for \(\alpha ,\beta ,\alpha +\beta \in \mathbb{R}\) .

Let \(\varphi\) be the automorphism of \(\mathfrak{g}\) considered in Def. 3 (iii). Then

\[
\begin{array}{r} \mathrm{N} _ {- \alpha , - \beta} X _ {- \alpha - \beta} = [ X _ {- \alpha}, X _ {- \beta} ] = [ \varphi (X _ {\alpha}), \varphi (X _ {\beta}) ] = \varphi ([ X _ {\alpha}, X _ {\beta} ]) \\ = \varphi (\mathrm{N} _ {\alpha , \beta} X _ {\alpha + \beta}) = \mathrm{N} _ {\alpha , \beta} X _ {- \alpha - \beta} \end{array}
\]

so \(N_{-\alpha,-\beta}=N_{\alpha,\beta}\) . Now \(N_{\alpha,\beta}=\pm(q+1)\) by (5).

PROPOSITION 8. Let \((X_{\alpha})_{\alpha \in \mathbb{R}}\) be a Chevalley system for \((\mathfrak{g},\mathfrak{h})\) . Let \(\mathrm{M}\) be a \(\mathbf{Z}\) -submodule of \(\mathfrak{h}\) containing the \(H_{\alpha}\) and contained in the group of weights of \(R^{\vee}\) . Let \(\mathfrak{g}_{\mathbf{Z}}\) be the Z-submodule of g generated by M and the \(X_{\alpha}\) . Then \(\mathfrak{g}_{\mathbf{Z}}\) is a Z-Lie subalgebra of g, and the canonical map from \(g_{Z} \otimes_{Z} k\) to g is an isomorphism.

If \(\alpha, \beta \in \mathbb{R}\) are such that \(\alpha + \beta \in \mathbb{R}\) , then \(\mathrm{N}_{\alpha, \beta} \in \mathbf{Z}\) (Prop. 7). On the other hand, if \(\alpha \in \mathbb{R}\) and \(h \in \mathbb{M}\) , then \(\alpha(h) \in \mathbf{Z}\) (Chap. VI, §1, no. 9). This proves that \(\mathfrak{g}_{\mathbf{Z}}\) is a \(\mathbf{Z}\) -Lie subalgebra of \(\mathfrak{g}\) . On the other hand, M is a free abelian group of rank dim \(\mathfrak{h}\) (Algebra, Chap. VII, §3, Th. 1), so \(\mathfrak{g}_{\mathbf{Z}}\) is a free abelian group of rank dim \(\mathfrak{g}\) ; this implies the last assertion.

\section{SUBALGEBRAS OF SPLIT SEMI-SIMPLE LIE ALGEBRAS}

In this paragraph, we denote by \((\mathfrak{g},\mathfrak{h})\) a split semi-simple Lie algebra, and by R its root system.

\subsection{SUBALGEBRAS STABLE UNDER ad h}

Lemma 1. Let \(\mathrm{V}\) be a vector subspace of \(\mathfrak{g}\) and \(\mathrm{R(V)}\) the set of \(\alpha \in \mathrm{R}\) such that \(\mathfrak{g}^{\alpha} \subset \mathrm{V}\) . Then, \((\mathrm{V} \cap \mathfrak{h}) + \sum_{\alpha \in \mathrm{R(V)}} \mathfrak{g}^{\alpha}\) is the largest vector subspace of \(\mathrm{V}\) stable under \(\operatorname{ad} \mathfrak{h}\) .

A vector subspace W of V is stable under ad h if and only if

\[
W = (W \cap \mathfrak {h}) + \sum_ {\alpha \in R} (W \cap \mathfrak {g} ^ {\alpha})
\]

(Algebra, Chap. VII, §2, no. 2, Cor. 1 of Th. 1). The largest vector subspace of V stable under ad \(\mathfrak{h}\) is thus \((\mathrm{V} \cap \mathfrak{h}) + \sum_{\alpha \in \mathrm{R}} (\mathrm{V} \cap \mathfrak{g}^{\alpha})\) . But \(\mathrm{V} \cap \mathfrak{g}^{\alpha} = \mathfrak{g}^{\alpha}\) for \(\alpha \in \mathrm{R}(\mathrm{V})\) , and \(\mathrm{V} \cap \mathfrak{g}^{\alpha} = 0\) for \(\alpha \notin \mathrm{R}(\mathrm{V})\) since \(\dim \mathfrak{g}^{\alpha} = 1\) . Q.E.D.

For any subset P of R, put

\[
\mathfrak {g} ^ {\mathrm{P}} = \sum_ {\alpha \in \mathrm{P}} \mathfrak {g} ^ {\alpha} \quad \mathfrak {h} _ {\mathrm{P}} = \sum_ {\alpha \in \mathrm{P}} \mathfrak {h} _ {\alpha}.
\]

If \(\mathrm{P}\subset \mathrm{R}\) and \(\mathbf{Q}\subset \mathbb{R}\) , we clearly have

\[
\begin{array}{c} {[ \mathfrak {h}, \mathfrak {g} ^ {\mathrm{P}} ] \subset \mathfrak {g} ^ {\mathrm{P}}} \\ {[ \mathfrak {g} ^ {\mathrm{P}}, \mathfrak {g} ^ {\mathrm{Q}} ] = \mathfrak {g} ^ {(\mathrm{P+Q}) \cap \mathrm{R}} + \mathfrak {h} _ {\mathrm{P} \cap (- \mathrm{Q})}.} \end{array}\tag{1}
\]

\begin{enumerate}
\def\labelenumi{(\arabic{enumi})}
\setcounter{enumi}{1}
\tightlist
\item
\end{enumerate}

Recall (Chap. VI, §1, no. 7, Def. 4) that a subset P of R is said to be closed if the conditions \(\alpha \in P\) , \(\beta \in P\) , \(\alpha + \beta \in R\) imply \(\alpha + \beta \in P\) , in other words if \((P + P) \cap R \subset P\) .

Lemma 2. Let \(\mathfrak{h}'\) be a vector subspace of \(\mathfrak{h}\) and \(\mathrm{P}\) a subset of \(\mathrm{R}\) . Then \(\mathfrak{h}' + \mathfrak{g}^{\mathrm{P}}\) is a subalgebra of \(\mathfrak{g}\) if and only if \(\mathrm{P}\) is a closed subset of \(\mathrm{R}\) and

\[
\mathfrak {h} ^ {\prime} \supset \mathfrak {h} _ {\mathrm{P} \cap (- \mathrm{P})}.
\]

Indeed,

\[
[ \mathfrak {h} ^ {\prime} + \mathfrak {g} ^ {\mathrm{P}}, \mathfrak {h} ^ {\prime} + \mathfrak {g} ^ {\mathrm{P}} ] = [ \mathfrak {h} ^ {\prime}, \mathfrak {g} ^ {\mathrm{P}} ] + [ \mathfrak {g} ^ {\mathrm{P}}, \mathfrak {g} ^ {\mathrm{P}} ] = \mathfrak {h} _ {\mathrm{P} \cap (- \mathrm{P})} + [ \mathfrak {h} ^ {\prime}, \mathfrak {g} ^ {\mathrm{P}} ] + \mathfrak {g} ^ {(\mathrm{P} + \mathrm{P}) \cap \mathrm{R}}.
\]

Hence \(\mathfrak{h}' + \mathfrak{g}^{\mathrm{P}}\) is a subalgebra of \(\mathfrak{g}\) if and only if

\[
\mathfrak {h} _ {\mathrm{P} \cap (- \mathrm{P})} \subset \mathfrak {h} ^ {\prime} \text { and } \mathfrak {g} ^ {(\mathrm{P} + \mathrm{P}) \cap \mathrm{R}} \subset \mathfrak {g} ^ {\mathrm{P}}
\]

which proves the lemma.

PROPOSITION 1. (i) The subalgebras of g stable under ad h are the vector subspaces of the form \(h' + g^{P}\) , where P is a closed subset of R and \(h'\) is a vector subspace of h containing \(\mathfrak{h}_{\mathrm{P}\cap(-\mathrm{P})}\) .

\begin{enumerate}
\def\labelenumi{(\roman{enumi})}
\setcounter{enumi}{1}
\tightlist
\item
  Let \(\mathfrak{h}'\) , \(\mathfrak{h}''\) be vector subspaces of \(\mathfrak{h}\) and \(\mathrm{P},\mathrm{Q}\) closed subsets of \(\mathrm{R}\) , with \(\mathfrak{h}'\supset \mathfrak{h}_{\mathrm{P}\cap (-\mathrm{P})}\) , \(\mathfrak{h}''\subset \mathfrak{h}'\) and \(\mathrm{Q}\subset \mathrm{P}\) . Then \(\mathfrak{h}'' + \mathfrak{g}^{\mathrm{Q}}\) is an ideal of \(\mathfrak{h}' + \mathfrak{g}^{\mathrm{P}}\) if and only if
\end{enumerate}

\[
(P + Q) \cap R \subset Q \text { and } \mathfrak {h} _ {P \cap (- Q)} \subset \mathfrak {h} ^ {\prime \prime} \subset \bigcap_ {\alpha \in P, \alpha \notin Q} \operatorname{Ker} \alpha .
\]

Assertion (i) follows immediately from Lemmas 1 and 2. Let \(\mathfrak{h}',\mathfrak{h}''\) , P, Q be as in (ii). Then

\[
[ \mathfrak {h} ^ {\prime} + \mathfrak {g} ^ {\mathrm{P}}, \mathfrak {h} ^ {\prime \prime} + \mathfrak {g} ^ {\mathrm{Q}} ] = \mathfrak {h} _ {\mathrm{P} \cap (- \mathrm{Q})} + [ \mathfrak {h} ^ {\prime}, \mathfrak {g} ^ {\mathrm{Q}} ] + [ \mathfrak {h} ^ {\prime \prime}, \mathfrak {g} ^ {\mathrm{P}} ] + \mathfrak {g} ^ {(\mathrm{P} + \mathrm{Q}) \cap \mathrm{R}}.
\]

Hence, \(\mathfrak{h}'' + \mathfrak{g}^{\mathrm{Q}}\) is an ideal of \(\mathfrak{h}' + \mathfrak{g}^{\mathrm{P}}\) if and only if

\[
\mathfrak {h} _ {\mathrm{P} \cap (- \mathrm{Q})} \subset \mathfrak {h} ^ {\prime \prime}, [ \mathfrak {h} ^ {\prime \prime}, \mathfrak {g} ^ {\mathrm{P}} ] \subset \mathfrak {g} ^ {\mathrm{Q}}, \mathfrak {g} ^ {(\mathrm{P} + \mathrm{Q}) \cap \mathrm{R}} \subset \mathfrak {g} ^ {\mathrm{Q}}.
\]

This implies (ii).

PROPOSITION 2. Let \(\mathfrak{a}\) be a subalgebra of \(\mathfrak{g}\) stable under ad \(\mathfrak{h}\) , and let \(\mathfrak{h}' \subset \mathfrak{h}\) , \(\mathrm{P} \subset \mathrm{R}\) be such that \(\mathfrak{a} = \mathfrak{h}' + \mathfrak{g}^{\mathrm{P}}\) .

\begin{enumerate}
\def\labelenumi{(\roman{enumi})}
\item
  Let \(\mathfrak{k}\) be the set of \(x\in \mathfrak{h}'\) such that \(\alpha (x) = 0\) for all \(\alpha \in \mathrm{P}\cap (-\mathrm{P})\) . The radical of \(\mathfrak{a}\) is \(\mathfrak{k} + \mathfrak{g}^{\mathrm{Q}}\) , where \(\mathrm{Q}\) is the set of \(\alpha \in \mathrm{P}\) such that \(-\alpha \notin \mathrm{P}\) . Moreover, \(\mathfrak{g}^{\mathrm{Q}}\) is a nilpotent ideal of \(\mathfrak{a}\) .
\item
  a is semi-simple if and only if P = -P and \(h' = h_{P}\) .
\item
  \(\mathfrak{a}\) is solvable if and only if \(\mathrm{P} \cap (-\mathrm{P}) = \emptyset\) . In that case \([\mathfrak{a}, \mathfrak{a}] = \mathfrak{g}^{\mathrm{S}}\) , where
\end{enumerate}

\[
\mathrm{S} = ((\mathrm{P} + \mathrm{P}) \cap \mathrm{R}) \cup \{\alpha \in \mathrm{P} | \alpha (\mathfrak {h} ^ {\prime}) \neq 0 \}.
\]

\begin{enumerate}
\def\labelenumi{(\roman{enumi})}
\setcounter{enumi}{3}
\item
  \(\mathfrak{a}\) is reductive in \(\mathfrak{g}\) if and only if \(\mathrm{P} = -\mathrm{P}\) .
\item
  \(\mathfrak{a}\) consists of nilpotent elements if and only if \(\mathfrak{h}' = 0\) . Then \(\mathrm{P} \cap (-\mathrm{P}) = \emptyset\) , and \(\mathfrak{a}\) is nilpotent.
\end{enumerate}

We prove (v). If \(\mathfrak{a}\) consists of nilpotent elements, \(\mathfrak{a}\) is clearly nilpotent, and \(\mathfrak{h}' = 0\) since the elements of \(\mathfrak{h}\) are semi-simple. Assume that \(\mathfrak{h}' = 0\) . By Prop. 1 (i), \(P \cap (-P) = \varnothing\) . By Chap. VI, §1, no. 7, Prop. 22, there exists a chamber C of R such that \(P \subset R_{+}(C)\) . Hence, there exists an integer \(n > 0\) with the following properties: if \(\alpha_{1}, \ldots, \alpha_{n} \in P\) and \(\beta \in R \cup \{0\}\) , then

\[
\alpha_ {1} + \dots + \alpha_ {n} + \beta \notin \mathrm{R} \cup \{0 \}.
\]

This implies that every element of \(\mathfrak{g}^{\mathrm{P}}\) is nilpotent, hence (v).

We prove (iii). If \(\mathrm{P} \cap (-\mathrm{P}) = \varnothing\) , \(\mathfrak{g}^{\mathrm{P}}\) is a subalgebra of \(\mathfrak{g}\) (Prop. 1 (i)), and is nilpotent by (v). Now

\[
[ \mathfrak {a}, \mathfrak {a} ] = [ \mathfrak {h} ^ {\prime}, \mathfrak {g} ^ {\mathrm{P}} ] + [ \mathfrak {g} ^ {\mathrm{P}}, \mathfrak {g} ^ {\mathrm{P}} ] = [ \mathfrak {h} ^ {\prime}, \mathfrak {g} ^ {\mathrm{P}} ] + \mathfrak {g} ^ {(\mathrm{P} + \mathrm{P}) \cap \mathrm{R}} \subset \mathfrak {g} ^ {\mathrm{P}},
\]

so \(\mathfrak{a}\) is solvable and \([\mathfrak{a},\mathfrak{a}]\) is given by the formula in the proposition. If \(\mathrm{P}\cap (-\mathrm{P})\neq \emptyset\) , let \(\alpha \in \mathrm{P}\) be such that \(-\alpha \in \mathrm{P}\) . Then \(\mathfrak{h}_{\alpha} + \mathfrak{g}^{\alpha} + \mathfrak{g}^{-\alpha}\) is a simple subalgebra of \(\mathfrak{a}\) so \(\mathfrak{a}\) is not solvable.

We prove (i). Since P is closed, \((\mathrm{P} + \mathrm{Q}) \cap \mathrm{R} \subset \mathrm{P}\) . If \(\alpha \in P, \beta \in Q\) and \(\alpha + \beta \in R\) , we cannot have \(\alpha + \beta \in -P\) , for, P being closed, this would imply that \(-\beta = -(\alpha + \beta) + \alpha \in P\) whereas \(\beta \in Q\) ; thus, \((\mathrm{P} + \mathrm{Q}) \cap \mathrm{R} \subset \mathrm{Q}\) . This proves that \(g^{Q}\) is an ideal of a, nilpotent by (v). We have \(\mathrm{P} \cap (-\mathrm{Q}) = \varnothing\) , and \(\mathrm{P} \cap (-\mathrm{P}) = \mathrm{P} \cap \complement Q\) , so \(\mathfrak{h}_{\mathrm{P} \cap (-\mathrm{Q})} \subset \mathfrak{k} \subset \bigcap_{\alpha \in \mathrm{P}, \alpha \notin Q} \operatorname{Ker} \alpha\) . By Prop. 1 (ii),

\(\mathfrak{k} + \mathfrak{g}^{\mathrm{Q}}\) is an ideal of \(\mathfrak{a}\) . Since \(\mathrm{Q} \cap (-\mathrm{Q}) = \emptyset\) , this ideal is solvable by (iii). It is therefore contained in the radical \(\mathfrak{r}\) of \(\mathfrak{a}\) . Since \(\mathfrak{r}\) is stable under every derivation of \(\mathfrak{a}\) , \(\mathfrak{r}\) is stable under ad \(\mathfrak{h}\) . Hence there exists a subset \(S\) of \(P\) such that \(\mathfrak{r} = (\mathfrak{r} \cap \mathfrak{h}) + \mathfrak{g}^{\mathrm{S}}\) . Suppose that \(\alpha \in S\) and that \(-\alpha \in P\) . Then \(\mathfrak{h}_{\alpha} = [\mathfrak{g}^{\alpha}, \mathfrak{g}^{-\alpha}] \subset \mathfrak{r}\) , so \(\mathfrak{g}^{-\alpha} = [\mathfrak{h}_{\alpha}, \mathfrak{g}^{-\alpha}] \subset \mathfrak{r} = 0\) , so that \(-\alpha \in S\) ; by (iii), this contradicts the fact that \(\mathfrak{r}\) is solvable. Consequently, \(S \subset Q\) . Finally, if \(x \in \mathfrak{r} \cap \mathfrak{h}\) and if \(\alpha \in P \cap (-P)\) , then \([x, \mathfrak{g}^{\alpha}] \subset \mathfrak{g}^{\alpha} \cap \mathfrak{r} = 0\) , so \(\alpha(x) = 0\) ; this shows that \(x \in \mathfrak{k}\) . Hence \(\mathfrak{r} \subset \mathfrak{k} + \mathfrak{g}^{\mathrm{Q}}\) and the proof of (i) is complete.

We prove (iv). By (i), the adjoint representation of \(\mathfrak{a}\) on \(\mathfrak{g}\) is semi-simple if and only if \(\mathrm{ad}_{\mathfrak{g}}x\) is semi-simple for all \(x\in \mathfrak{k} + \mathfrak{g}^{\mathrm{Q}}\) (Chap. I, §6, no. 5, Th. 4); by (v), this is the case if and only if \(\mathrm{Q} = \varnothing\) , in other words \(\mathrm{P} = -\mathrm{P}\) .

We prove (ii). If \(\mathfrak{a}\) is semi-simple, \(P = -P\) by (i), so \(\mathfrak{h}_{\mathrm{P}} \subset \mathfrak{h}'\) ; further, \(\mathfrak{a} = [\mathfrak{a}, \mathfrak{a}] \subset \mathfrak{h}_{\mathrm{P}} + \mathfrak{g}^{\mathrm{P}}\) and consequently \(\mathfrak{h}' = \mathfrak{h}_{\mathrm{P}}\) . If \(P = -P\) and \(\mathfrak{h}' = \mathfrak{h}_{\mathrm{P}}\) , \(\mathfrak{a}\) is reductive by (iv), and \(\mathfrak{a} = \sum_{\alpha \in P} \mathfrak{s}_{\alpha}\) , so \(\mathfrak{a} = [\mathfrak{a}, \mathfrak{a}]\) and \(\mathfrak{a}\) is semi-simple.

PROPOSITION 3. Let \(\mathfrak{a}\) be a semi-simple subalgebra of \(\mathfrak{g}\) stable under \(\operatorname{ad}(\mathfrak{h})\) and let \(\mathrm{P}\) be the subset of \(\mathrm{R}\) such that \(\mathfrak{a} = \mathfrak{h}_{\mathrm{P}} + \mathfrak{g}^{\mathrm{P}}\) .

\begin{enumerate}
\def\labelenumi{(\roman{enumi})}
\item
  \(\mathfrak{h}_{\mathrm{P}}\) is a splitting Cartan subalgebra of \(\mathfrak{a}\) .
\item
  The root system of \((\mathfrak{a},\mathfrak{h}_{\mathrm{P}})\) is the set of restrictions to \(\mathfrak{h}_{\mathrm{P}}\) of elements of \(\mathrm{P}\) .
\end{enumerate}

Since \(h_{P}\) is stable under ad h, its normalizer in a is stable under ad h, and hence is of the form \(h_{P} + g^{Q}\) where \(Q \subset P\) (Lemma 1). If \(\alpha \in Q\) ,

\[
\mathfrak {g} ^ {\alpha} = \left[ \mathfrak {h} _ {\alpha}, \mathfrak {g} ^ {\alpha} \right] \subset \left[ \mathfrak {h} _ {\mathrm{P}}, \mathfrak {g} ^ {\alpha} \right] \subset \mathfrak {h} _ {\mathrm{P}},
\]

which is absurd. Thus \(Q = \varnothing\) and \(h_{P}\) is its own normalizer in a. This proves that \(h_{P}\) is a Cartan subalgebra of a. If \(x \in h_{P}\) , \(ad_{g}x\) , and a fortiori \(ad_{a}x\) , are triangularizable. Thus (i) is proved, and (ii) is clear.

By Prop. 1 (i), the subalgebras of g containing h are the sets \(h + g^{P}\) where P is a closed subset of R. By Chap. VII, §3, Prop. 3, every Cartan subalgebra of \(h + g^{P}\) is a Cartan subalgebra of g.

PROPOSITION 4. Let \(\mathfrak{a}\) be a subalgebra of \(\mathfrak{g}\) containing \(\mathfrak{h}\) , \(x\) an element of \(\mathfrak{a}\) , \(s\) and \(n\) its semi-simple and nilpotent components. Then \(s \in \mathfrak{a}\) and \(n \in \mathfrak{a}\) .

We have \((\operatorname{ad} x)\mathfrak{a} \subset \mathfrak{a}\) , so \((\operatorname{ad} s)\mathfrak{a} \subset \mathfrak{a}\) and \((\operatorname{ad} n)\mathfrak{a} \subset \mathfrak{a}\) . Since a is its own normalizer in g (Chap. VII, §2, no. 1, Cor. 4 of Prop. 4), \(s \in a\) and \(n \in a\) .

PROPOSITION 5. Let P be a closed subset of R.

\begin{enumerate}
\def\labelenumi{(\roman{enumi})}
\item
  \(\mathfrak{h} + \mathfrak{g}^{\mathrm{P}}\) is solvable if and only if \(\mathrm{P} \cap (-\mathrm{P}) = \varnothing\) . In that case, \([\mathfrak{h} + \mathfrak{g}^{\mathrm{P}}, \mathfrak{h} + \mathfrak{g}^{\mathrm{P}}] = \mathfrak{g}^{\mathrm{P}}\) .
\item
  \(\mathfrak{h} + \mathfrak{g}^{\mathrm{P}}\) is reductive if and only if \(\mathrm{P} = -\mathrm{P}\) .
\end{enumerate}

Assertion (i) follows from Prop. 2 (iii). If \(\mathrm{P} = -\mathrm{P}\) , \(\mathfrak{h} + \mathfrak{g}^{\mathrm{P}}\) is reductive (Prop. 2 (iv)). Assume that \(\mathfrak{a} = \mathfrak{h} + \mathfrak{g}^{\mathrm{P}}\) is reductive. Then

\[
\mathfrak {g} ^ {\mathrm{P}} = \left[ \mathfrak {h}, \mathfrak {g} ^ {\mathrm{P}} \right] \subset \left[ \mathfrak {a}, \mathfrak {a} \right] \subset \mathfrak {h} + \mathfrak {g} ^ {\mathrm{P}},
\]

so \([a, a]\) is of the form \(h' + g^{P}\) with \(h' \subset h\) ; since \([a, a]\) is semi-simple, P = -P (Prop. 2 (ii)).

\subsection{IDEALS}

PROPOSITION 6. Let \(\mathrm{R}_1, \ldots, \mathrm{R}_p\) be the irreducible components of \(\mathbb{R}\) . For \(i = 1, \ldots, p\) , put \(\mathfrak{g}_i = \mathfrak{h}_{\mathrm{R}_i} + \mathfrak{g}^{\mathrm{R}_i}\) . Then \(\mathfrak{g}_1, \ldots, \mathfrak{g}_p\) are the simple components of \(\mathfrak{g}\) .

The \(g_{i}\) are ideals of g (Prop. 1 (ii)). It is clear that g is the direct sum of the \(g_{i}\) , hence the product of the \(g_{i}\) . Let a and b be complementary ideals of g. Then a and b are semi-simple and stable under ad h, so there exist subsets P, Q of R such that \(a = h_{P} + g^{P}\) , \(b = h_{Q} + g^{Q}\) . Then \(h_{P}\) , \(h_{Q}\) are orthogonal complements of each other in h for the Killing form, so P and Q are unions of irreducible components of R. This proves that the \(g_{i}\) are minimal ideals of g.

COROLLARY 1. g is simple if and only if R is irreducible (in other words, its Dynkin graph is connected).

This follows from Prop. 6.

A Lie algebra \(\mathfrak{a}\) is said to be absolutely simple if, for every extension \(k'\) of \(k\) , the \(k'\) -Lie algebra \(\mathfrak{a}_{(k')}\) is simple.

COROLLARY 2. A splittable simple Lie algebra is absolutely simple.

This follows from Cor. 1.

If g is of type \(A_{l}\) ( \(l \geq 1\) ) or \(B_{l}\) ( \(l \geq 1\) ) or \(C_{l}\) ( \(l \geq 1\) ) or \(D_{l}\) ( \(l \geq 3\) ), g is said to be a classical splittable simple Lie algebra. If g is of type \(E_{6}\) , \(E_{7}\) , \(E_{8}\) , \(F_{4}\) , or \(G_{2}\) , g is said to be an exceptional splittable simple Lie algebra.

\subsection{BOREL SUBALGEBRAS}

PROPOSITION 7. Let \(\mathfrak{b} = \mathfrak{h} + \mathfrak{g}^{\mathrm{P}}\) be a subalgebra of \(\mathfrak{g}\) containing \(\mathfrak{h}\) . The following conditions are equivalent:

\begin{enumerate}
\def\labelenumi{(\roman{enumi})}
\item
  \(\mathfrak{b}\) is a maximal solvable subalgebra of \(\mathfrak{g}\) ;
\item
  there exists a chamber C of R such that \(P = R_{+}(C)\) ;
\item
  \(\mathrm{P} \cap (-\mathrm{P}) = \varnothing\) and \(\mathrm{P} \cup (-\mathrm{P}) = \mathrm{R}\) .
\item
  \(\Longrightarrow\) (ii): If \(\mathfrak{b}\) is solvable, \(P \cap (-P) = \varnothing\) . Then there exists a chamber \(C\) of \(R\) such that \(P \subset R_{+}(C)\) (Chap. VI, §1, no. 7, Prop. 22). Then \(\mathfrak{h} + \mathfrak{g}^{R_{+}(C)}\) is a solvable subalgebra of \(\mathfrak{g}\) containing \(\mathfrak{b}\) , hence equal to \(\mathfrak{b}\) if \(\mathfrak{b}\) is maximal.
\item
  \(\Longrightarrow\) (iii): This is clear.
\item
  \(\Longrightarrow\) (i): Assume that \(P \cap (-P) = \varnothing\) and that \(P \cup (-P) = R\) . Then b is solvable. Let \(b'\) be a solvable subalgebra of g containing b. There exists a subset Q of R such that \(b' = h + g^{Q}\) . Then \(Q \cap (-Q) = \varnothing\) and \(Q \supset P\) , so Q = P and \(b' = b\) .
\end{enumerate}

DEFINITION 1. A subalgebra of g containing h and satisfying the equivalent condition in Prop. 7 is called a Borel subalgebra of (g, h).

A subalgebra \(\mathfrak{b}\) of a splittable algebra \(\mathfrak{g}\) is called a Borel subalgebra of \(\mathfrak{g}\) if there exists a splitting Cartan subalgebra \(\mathfrak{h}'\) of \(\mathfrak{g}\) such that \(\mathfrak{b}\) is a Borel subalgebra of \((\mathfrak{g},\mathfrak{h}')\) .

Let \((\mathfrak{g},\mathfrak{h})\) be a split reductive Lie algebra. Let \(\mathfrak{g}=\mathfrak{c}\times\mathfrak{s}\) with c commutative and s semi-simple. A subalgebra of g of the form \(\mathfrak{c}\times\mathfrak{b}\) , where b is a Borel subalgebra of \((\mathfrak{s},\mathfrak{h}\cap\mathfrak{s})\) , is called a Borel subalgebra of \((\mathfrak{g},\mathfrak{h})\) .

With the notations of Prop. 7, we also say that b is the Borel subalgebra of g defined by h and C (or by h and the basis of R associated to C).

Remark. The map which associates \(\mathrm{R}_{+}(\mathrm{C})\) to a chamber C of R is injective (Chap. VI, §1, no. 7, Cor. 1 of Prop. 20). Consequently, \(\mathrm{C} \mapsto \mathfrak{h} + \mathfrak{g}^{\mathrm{R}_{+}(\mathrm{C})}\) is a bijection from the set of chambers of R to the set of Borel subalgebras of \((\mathfrak{g}, \mathfrak{h})\) . Thus, the number of Borel subalgebras of \((\mathfrak{g}, \mathfrak{h})\) is equal to the order of the Weyl group of R (Chap. VI, §1, no. 5, Th. 2).

PROPOSITION 8. Let \(\mathfrak{b}\) be a subalgebra of \(\mathfrak{g}\) , \(k'\) an extension of \(k\) . Then \(\mathfrak{b} \otimes_k k'\) is a Borel subalgebra of \((\mathfrak{g} \otimes_k k', \mathfrak{h} \otimes_k k')\) if and only if \(\mathfrak{b}\) is a Borel subalgebra of \((\mathfrak{g}, \mathfrak{h})\) .

This is clear from condition (iii) of Prop. 7.

PROPOSITION 9. Let \(\mathfrak{b}\) be the Borel subalgebra of \((\mathfrak{g},\mathfrak{h})\) defined by a chamber C of R. Let \(\mathfrak{n} = \mathfrak{g}^{\mathrm{R}_{+}(\mathrm{C})} = \sum_{\alpha \in \mathrm{R},\alpha >0}\mathfrak{g}^{\alpha}\) . Let \(l = \dim \mathfrak{h}\) .

\begin{enumerate}
\def\labelenumi{(\roman{enumi})}
\item
  If \(h \in \mathfrak{h}\) and \(x \in \mathfrak{n}\) , the characteristic polynomial of \(\mathrm{ad}_{\mathfrak{g}}(h + x)\) is \(\mathrm{T}^l \prod_{\alpha \in \mathbb{R}} (\mathrm{T} - \alpha(h))\) .
\item
  The largest nilpotent ideal of \(\mathfrak{b}\) is equal to \(\mathfrak{n}\) and to \([\mathfrak{b},\mathfrak{b}]\) . This is also the set of elements of \(\mathfrak{b}\) nilpotent in \(\mathfrak{g}\) .
\item
  Let \(\mathrm{B}\) be the basis of \(\mathbb{R}\) associated to \(\mathbb{C}\) . For all \(\alpha \in \mathbb{B}\) , let \(X_{\alpha}\) be a non-zero element of \(\mathfrak{g}^{\alpha}\) . Then \((X_{\alpha})_{\alpha \in \mathbb{B}}\) generates the Lie algebra \(\mathfrak{n}\) . We have \([\mathfrak{n}, \mathfrak{n}] = \sum_{\alpha \in \mathbb{R}, \alpha > 0, \alpha \notin \mathbb{B}} \mathfrak{g}^{\alpha}\) .
\end{enumerate}

There exists a total order on \(\mathfrak{h}_{\mathbf{Q}}^{*}\) compatible with its vector space structure and such that the elements of \(\mathrm{R}_{+}(\mathrm{C})\) are \(>0\) (Chap. VI, §1, no. 7). Let \(h, x\) be as in (i) and \(y \in \mathfrak{g}^{\alpha}\) . Then \([h + x, y] = \alpha(h)y + z\) where \(z \in \sum_{\beta > \alpha} \mathfrak{g}^{\beta}\) . Then, with respect to a suitable basis of \(\mathfrak{g}\) , the matrix of \(\operatorname{ad}_{\mathfrak{g}}(h + x)\) has the following properties:

\begin{enumerate}
\def\labelenumi{\arabic{enumi})}
\item
  it is lower triangular;
\item
  the diagonal entries of the matrix are the number 0 (l times) and the \(\alpha(h)\) for \(\alpha \in \mathbb{R}\) .
\end{enumerate}

This proves (i). It also shows that the characteristic polynomial of \(\operatorname{ad}_{\mathfrak{b}}(h+x)\) is \(\mathrm{T}^{l}\prod_{\alpha\in\mathrm{R}_{+}(\mathrm{C})}(\mathrm{T}-\alpha(h))\) . It follows from the preceding that the set of elements of b nilpotent in g, as well as the largest nilpotent ideal of b, are equal to n.~We have \(n=[b,b]\) by Prop. 5 (i). Finally, assertion (iii) follows from §2, Prop. 4 (ii) and Chap. VI, §1, no. 6, Prop. 19.

COROLLARY. Let \(\mathfrak{b}\) be a Borel subalgebra of \(\mathfrak{g}\) .

\begin{enumerate}
\def\labelenumi{(\roman{enumi})}
\item
  Every Cartan subalgebra of b is a splitting Cartan subalgebra of g.
\item
  If \(\mathfrak{h}_1, \mathfrak{h}_2\) are Cartan subalgebras of \(\mathfrak{b}\) , there exists \(x \in [\mathfrak{b}, \mathfrak{b}]\) such that \(e^{\mathrm{ad}_{\mathfrak{g}} x} \mathfrak{h}_1 = \mathfrak{h}_2\) .
\end{enumerate}

Assertion (i) follows from Prop. 9 (i) and Chap. VII, §3, no. 3, Prop. 3. Assertion (ii) follows from Prop. 9 (ii) and Chap. VII, §3, no. 4, Th. 3.

PROPOSITION 10. Let \(\mathfrak{b},\mathfrak{b}'\) be Borel subalgebras of \(\mathfrak{g}\) . There exists a splitting Cartan subalgebra of \(\mathfrak{g}\) contained in \(\mathfrak{b} \cap \mathfrak{b}'\) .

Let \(\mathfrak{h}\) be a Cartan subalgebra of \(\mathfrak{b}\) , \(\mathfrak{n} = [\mathfrak{b},\mathfrak{b}],\mathfrak{n}' = [\mathfrak{b}',\mathfrak{b}']\) , \(\mathfrak{p} = \mathfrak{b}\cap \mathfrak{b}'\) , and \(\mathfrak{s}\) a vector subspace of \(\mathfrak{g}\) complementary to \(\mathfrak{b} + \mathfrak{b}'\) . Denote by \(\mathfrak{s}^{\perp},\mathfrak{b}^{\perp},\mathfrak{b}'^{\perp}\) the orthogonal complements of \(\mathfrak{s},\mathfrak{b},\mathfrak{b}^{\prime}\) with respect to the Killing form of \(\mathfrak{g}\) . Put \(l = \dim \mathfrak{h}, n = \dim \mathfrak{n}, p = \dim \mathfrak{p}\) . Then \(\dim \mathfrak{b} = \dim \mathfrak{b}^{\prime} = l + n\) ,

\[
\dim \mathfrak {s} ^ {\perp} = \dim (\mathfrak {b} + \mathfrak {b} ^ {\prime}) = 2 (l + n) - p,
\]

and so

\[
\begin{array}{c} \dim (\mathfrak {s} ^ {\perp} \cap \mathfrak {p}) \geq \dim \mathfrak {s} ^ {\perp} + \dim \mathfrak {p} - \dim \mathfrak {g} \\ = 2 (l + n) - p + p - (l + 2 n) = l. \end{array}\tag{3}
\]

By Prop. 1 of §2, no. 2, \(\mathfrak{n} \subset \mathfrak{b}^{\perp}\) , \(\mathfrak{n}' \subset \mathfrak{b}'^{\perp}\) . The elements of \(\mathfrak{p} \cap \mathfrak{n}\) are nilpotent in \(\mathfrak{g}\) (Prop. 9 (ii)), and belong to \(\mathfrak{b}'\) , and hence to \(\mathfrak{n}'\) (Prop. 9 (ii)). Consequently, \(\mathfrak{p} \cap \mathfrak{n} \subset \mathfrak{n} \cap \mathfrak{n}' \subset \mathfrak{b}^{\perp} \cap \mathfrak{b}'^{\perp}\) , so \(\mathfrak{s}^{\perp} \cap \mathfrak{p} \cap \mathfrak{n} = 0\) . In view of (3), we see that \(\mathfrak{s}^{\perp} \cap \mathfrak{p}\) is a complement of \(\mathfrak{n}\) in \(\mathfrak{b}\) . Let \(z\) be an element of \(\mathfrak{h}\) regular in \(\mathfrak{g}\) ; there exists \(y \in \mathfrak{n}\) such that \(y + z \in \mathfrak{s}^{\perp} \cap \mathfrak{p}\) ; by Prop. 9 (i), \(\operatorname{ad}_{\mathfrak{g}}(y + z)\) has the same characteristic polynomial as \(\operatorname{ad}_{\mathfrak{g}}z\) , so \(x = y + z\) is regular in \(\mathfrak{g}\) and a fortiori in \(\mathfrak{b}\) and \(\mathfrak{b}'\) (Chap. VII, §2, no. 2, Prop. 9). Since \(\mathfrak{g}, \mathfrak{b}, \mathfrak{b}'\) have the same rank, \(\mathfrak{b}^{0}(x) = \mathfrak{g}^{0}(x) = \mathfrak{b}'^{0}(x)\) is simultaneously a Cartan subalgebra of \(\mathfrak{b}\) , of \(\mathfrak{g}\) and of \(\mathfrak{b}'\) (Chap. VII, §3, no. 3, Th. 2). Finally, this Cartan subalgebra of \(\mathfrak{g}\) is splitting by the Cor. of Prop. 9.

COROLLARY. The group \(\operatorname{Aut}_{e}(\mathfrak{g})\) operates transitively on the set of pairs \((\mathfrak{t},\mathfrak{b})\) where t is a splitting Cartan subalgebra of g and b is a Borel subalgebra of \((\mathfrak{g},\mathfrak{t})\) .

Let \((\mathfrak{t}_1, \mathfrak{b}_1)\) and \((\mathfrak{t}_2, \mathfrak{b}_2)\) be two such pairs. There exists a splitting Cartan subalgebra \(\mathfrak{t}\) of \(\mathfrak{g}\) contained in \(\mathfrak{b}_1 \cap \mathfrak{b}_2\) (Prop. 10). By the Cor. of Prop. 9, we are reduced to the case in which \(\mathfrak{t}_1 = \mathfrak{t}_2 = \mathfrak{t}\) . Let \(S\) be the root system of \((\mathfrak{g}, \mathfrak{t})\) . There exists bases \(B_1, B_2\) of \(S\) such that \(\mathfrak{b}_i\) is associated to \(B_i\) ( \(i = 1, 2\) ), and there exists \(s \in W(S)\) which transforms \(B_1\) into \(B_2\) . Finally, there exists \(a \in \operatorname{Aut}_e(\mathfrak{g})\) such that \(a|t = s\) ( \(\S 2\) , no. 2, Cor. of Th. 2). Then \(a(t) = t\) and \(a(\mathfrak{b}_1) = \mathfrak{b}_2\) .

\subsection{PARABOLIC SUBALGEBRAS}

PROPOSITION 11. Let \(\mathfrak{p} = \mathfrak{h} + \mathfrak{g}^{\mathrm{P}}\) be a subalgebra of \(\mathfrak{g}\) containing \(\mathfrak{h}\) . The following conditions are equivalent:

\begin{enumerate}
\def\labelenumi{(\roman{enumi})}
\item
  p contains a Borel subalgebra of \((\mathfrak{g}, \mathfrak{h})\) ;
\item
  there exists a chamber C of R such that \(\mathrm{P} \supset \mathrm{R}_{+}(\mathrm{C})\) ;
\item
  P is parabolic, in other words (Chap. VI, §1, no. 7, Def. 4), \(\mathrm{P} \cup (-\mathrm{P}) = \mathrm{R}\) .
\end{enumerate}

Conditions (i) and (ii) are equivalent by Prop. 7. Conditions (ii) and (iii) are equivalent by Chap. VI, §1, no. 7, Prop. 20.

DEFINITION 2. A subalgebra of g containing h and satisfying the equivalent conditions of Prop. 11 is called a parabolic subalgebra of \((\mathfrak{g}, \mathfrak{h})\) . A parabolic subalgebra of g is a parabolic subalgebra of \((\mathfrak{g}, \mathfrak{h}')\) where \(h'\) is a splitting Cartan subalgebra of g.

This definition extends immediately to the case in which \((\mathfrak{g}, \mathfrak{h})\) is a split reductive Lie algebra.

Remark. Let B be a basis of R, and b the corresponding Borel subalgebra. If \(\Sigma \subset B\) , denote by \(Q_{\Sigma}\) the set of roots that are linear combinations of elements of \(\Sigma\) with coefficients \(\leq 0\) ; put \(\mathfrak{p}(\Sigma) = \mathrm{R}_{+}(\mathrm{B}) \cup \mathrm{Q}_{\Sigma}\) and \(\mathfrak{p}_{\Sigma} = \mathfrak{h} \oplus \mathfrak{g}^{\mathrm{P}(\Sigma)}\) . By Chap. VI, §1, no. 7, Lemma 3 and Prop. 20, \(p_{\Sigma}\) is a parabolic subalgebra containing b and every parabolic subalgebra of g containing b is obtained in this way.

Lemma 3. Let V be a finite dimensional real vector space, S a root system in \(V^{*}\) , P the set of parabolic subsets of S; let H be the set of Ker \(\alpha\) for \(\alpha \in S\) , and F the set of facets of V relative to H (Chap. V, §1, no. 2, Def. 1).

If \(\mathrm{P} \in \mathcal{P}\) , let \(\overline{\mathrm{F}}(\mathrm{P})\) be the set of \(v \in \mathrm{V}\) such that \(\alpha(v) \geq 0\) for all \(\alpha \in \mathrm{P}\) . If \(\mathrm{F} \in \mathcal{F}\) , let \(\mathrm{P}(\mathrm{F})\) be the set of \(\alpha \in \mathrm{R}\) such that \(\alpha(v) \geq 0\) for all \(v \in \mathrm{F}\) .

Then \(\mathrm{F} \mapsto \mathrm{P}(\mathrm{F})\) is a bijection from \(\mathcal{F}\) to \(\mathcal{P}\) ; for all \(\mathrm{F} \in \mathcal{F}\) , \(\overline{\mathrm{F}}(\mathrm{P}(\mathrm{F}))\) is the closure of \(\mathrm{F}\) .

\begin{enumerate}
\def\labelenumi{\alph{enumi})}
\tightlist
\item
  Let \(P \in P\) . There exists a chamber C of S and a subset \(\Sigma\) of the basis \(\mathrm{B}(\mathrm{C})\) such that \(\mathrm{P} = \mathrm{S}_{+}(\mathrm{C}) \cup \mathrm{Q}\) where Q is the set of linear combinations of elements of \(\Sigma\) with non-positive integer coefficients (Chap. VI, §1, no. 7, Prop. 20). Put
\end{enumerate}

\[
\mathrm{B} (\mathrm{C}) = \{\alpha_ {1}, \dots , \alpha_ {l} \}, \Sigma = \{\alpha_ {1}, \dots , \alpha_ {m} \}.
\]

If \(v \in \mathrm{V}\) , we have the following equivalences:

\[
\begin{array}{c} \alpha (v) \geq 0 \text {   for   all   } \alpha \in \mathrm{P} \\ \Longleftrightarrow \alpha_ {1} (v) \geq 0, \ldots \alpha_ {l} (v) \geq 0, \alpha_ {1} (v) \leq 0, \ldots , \alpha_ {m} (v) \leq 0 \\ \Longleftrightarrow \alpha_ {1} (v) = \dots = \alpha_ {m} (v) = 0, \alpha_ {m + 1} (v) \geq 0, \ldots , \alpha_ {l} (v) \geq 0, \end{array}
\]

so \(\overline{\mathrm{F}}(\mathrm{P})\) is the closure of the set

\[
\{v \in \mathrm{V} \mid \alpha_ {1} (v) = \dots = \alpha_ {m} (v) = 0, \alpha_ {m + 1} (v) > 0, \dots , \alpha_ {l} (v) > 0 \},
\]

a set which is a facet F relative to H since every element of S is a linear combination of \(\alpha_{1},\ldots,\alpha_{l}\) in which the coefficients are either all \(\geq0\) or all \(\leq0\) . Moreover, if \(\beta=u_{1}\alpha_{1}+\cdots+u_{l}\alpha_{l}\in S\) ,

\[
\begin{array}{c} \beta \in \mathrm{P} (\mathrm{F}) \Longleftrightarrow u _ {m + 1} \geq 0, \ldots , u _ {l} \geq 0 \\ \Longleftrightarrow \beta \in \mathrm{S} _ {+} (\mathrm{C}) \text {   or   } (- \beta \in \mathrm{S} _ {+} (\mathrm{C}) \text {   and   } u _ {m + 1} = \ldots = u _ {l} = 0) \\ \Longleftrightarrow \beta \in \mathrm{S} _ {+} (\mathrm{C}) \cup \mathrm{Q} = \mathrm{P}, \end{array}
\]

so P(F) = P.

\begin{enumerate}
\def\labelenumi{\alph{enumi})}
\setcounter{enumi}{1}
\tightlist
\item
  Let \(F \in F\) . It is clear that \(P(F) \in \mathcal{P}\) . On the other hand, F is contained in the closure of a chamber relative to \(\mathcal{H}(\text{Chap. V}, \S1, \text{no. 3, formulas (6)}),\) and so is a facet relative to the set of walls of this chamber (Chap. V, \(\S1\) , no. 4, Prop. 9). Consequently, \(\overline{F}\) is of the form \(\{v \in V \mid \alpha(v) \geq 0 \text{ for all } \alpha \in T\}\) , where T is a subset of S which we can clearly take to be equal to \(P(F)\) . Thus, \(\overline{F} = \overline{F}(P(F))\) . Q.E.D.
\end{enumerate}

If \(P \in P\) , the facet F such that \(P = P(F)\) is said to be associated to P; we denote it by \(F(P)\) . We extend these conventions to the case in which \((\mathfrak{g}, \mathfrak{h})\) is split reductive.

PROPOSITION 12. Let \(\mathcal{H}\) be the set of hyperplanes of \(\mathfrak{h}_{\mathbf{R}}\) consisting of the kernels of the roots in R. Let \(\mathcal{F}\) be the set of facets of \(\mathfrak{h}_{\mathbf{R}}\) relative to \(\mathcal{H}\) . Let \(\mathcal{S}\) be the set of parabolic subalgebras of \((\mathfrak{g},\mathfrak{h})\) . For every \(\mathfrak{p} = \mathfrak{h} + \mathfrak{g}^{\mathrm{P}}\in \mathcal{S}\) , let \(\mathrm{F}(\mathfrak{p})\) be the facet associated to P. Then \(\mathfrak{p}\mapsto \mathrm{F}(\mathfrak{p})\) is a bijection from \(\mathcal{S}\) to \(\mathcal{F}\) . If \(\mathfrak{p}_1,\mathfrak{p}_2\in \mathcal{P}\) ,

\[
\mathfrak {p} _ {1} \supset \mathfrak {p} _ {2} \Longleftrightarrow \mathrm{F} (\mathfrak {p} _ {1}) \subset \overline {{\mathrm{F} (\mathfrak {p} _ {2})}}.
\]

This follows immediately from Lemma 3.

Example. The facets corresponding to the parabolic subalgebras of \((\mathfrak{g}, \mathfrak{h})\) containing a Borel algebra b are the facets contained in the closure of the chamber associated to b (cf.~the Remark above).

PROPOSITION 13. Let \(\mathfrak{p} = \mathfrak{h} + \mathfrak{g}^{\mathrm{P}}\) be a parabolic subalgebra of \((\mathfrak{g},\mathfrak{h})\) , Q the set of \(\alpha \in \mathrm{P}\) such that \(-\alpha \notin \mathrm{P}\) , and \(\mathfrak{s} = \mathfrak{h} + \mathfrak{g}^{\mathrm{P}\cap (-\mathrm{P})}\) . Then \(\mathfrak{p} = \mathfrak{s} \oplus \mathfrak{g}^{\mathrm{Q}}\) , \(\mathfrak{s}\) is reductive in \(\mathfrak{g}\) , and \(\mathfrak{g}^{\mathrm{Q}}\) is the largest nilpotent ideal of \(\mathfrak{p}\) and the nilpotent radical of \(\mathfrak{p}\) . The centre of \(\mathfrak{p}\) is zero.

By Prop. 2, \(\mathfrak{s}\) is reductive in \(\mathfrak{g}\) and \(\mathfrak{g}^{\mathrm{Q}}\) is a nilpotent ideal of \(\mathfrak{p}\) . If \(\mathfrak{n}\) is the largest nilpotent ideal of \(\mathfrak{p}\) , \(\mathfrak{g}^{\mathrm{Q}} \subset \mathfrak{n} \subset \mathfrak{h} + \mathfrak{g}^{\mathrm{Q}}\) (Prop. 2 (i)); if \(x \in \mathfrak{n} \cap \mathfrak{h}\) , \(\operatorname{ad}_{\mathfrak{p}} x\) is nilpotent, so \(\alpha(x) = 0\) for all \(\alpha \in \mathrm{P}\) , and hence \(x = 0\) ; this proves that \(\mathfrak{n} = \mathfrak{g}^{\mathrm{Q}}\) . Since \([\mathfrak{h}, \mathfrak{g}^{\mathrm{Q}}] = \mathfrak{g}^{\mathrm{Q}}\) , the nilpotent radical of \(\mathfrak{p}\) contains \(\mathfrak{g}^{\mathrm{Q}}\) and consequently is equal to \(\mathfrak{g}^{\mathrm{Q}}\) . Let \(z = h + \sum_{\alpha \in \mathrm{P}} u_{\alpha}\) (where \(h \in \mathfrak{h}, u_{\alpha} \in \mathfrak{g}^{\alpha}\) ) be an element of the centre of \(\mathfrak{p}\) . For all \(h' \in \mathfrak{h}\) , \(0 = [h', z] = \sum \alpha(h')u_{\alpha}\) , so \(u_{\alpha} = 0\) for all \(\alpha \in \mathrm{P}\) ; it follows that \([h, \mathfrak{g}^{\beta}] = 0\) for all \(\beta \in \mathrm{P}\) , so \(h = 0\) .

\subsection{NON-SPLIT CASE}

PROPOSITION 14. Let \(k'\) be an extension of \(k\) and \(\mathfrak{g}' = \mathfrak{g} \otimes_k k'\) . Let \(\mathfrak{m}\) be a subalgebra of \(\mathfrak{g}\) and \(\mathfrak{m}' = \mathfrak{m} \otimes_k k'\) . If \(\mathfrak{m}'\) is a parabolic (resp. Borel) subalgebra of \(\mathfrak{g}'\) , \(\mathfrak{m}\) is a parabolic (resp. Borel) subalgebra of \(\mathfrak{g}\) .

By Prop. 8 and 11, it suffices to prove that m contains a splitting Cartan subalgebra of g. Let b be a Borel subalgebra of g. Then \(b' = b \otimes_{k} k'\) is a Borel subalgebra of \(\mathfrak{g}'\) , so \(\mathfrak{m}' \cap \mathfrak{b}'\) contains a Cartan subalgebra of \(\mathfrak{g}'\) (Prop. 10). Let \(\mathfrak{t}\) be a Cartan subalgebra of \(\mathfrak{m} \cap \mathfrak{b}\) . Then \(\mathfrak{t} \otimes_k k'\) is a Cartan subalgebra of \(\mathfrak{m}' \cap \mathfrak{b}'\) , and hence of \(\mathfrak{g}'\) (Chap. VII, §3, no. 3, Prop. 3). Consequently, \(\mathfrak{t}\) is a Cartan subalgebra of \(\mathfrak{g}\) , and it is splitting since it is contained in \(\mathfrak{b}\) .

DEFINITION 3. Let \(\mathfrak{a}\) be a semi-simple (or more generally reductive) Lie algebra and \(\bar{k}\) an algebraic closure of \(k\) . A subalgebra \(\mathfrak{m}\) of \(\mathfrak{a}\) is said to be parabolic (resp. Borel) if \(\mathfrak{m} \otimes_k \bar{k}\) is a parabolic (resp. Borel) subalgebra of \(\mathfrak{a} \otimes_k \bar{k}\) .

If \(\mathfrak{a}\) is splittable, Prop. 14 shows that this definition is equivalent to Definition 2 (resp. to Definition 1).

PROPOSITION 15. Let \(\mathfrak{a}\) be a reductive Lie algebra, \(k'\) an extension of \(k\) , and \(\mathfrak{m}\) a subalgebra of \(\mathfrak{a}\) . Then \(\mathfrak{m}\) is a parabolic (resp. Borel) subalgebra of \(\mathfrak{a}\) if and only if \(\mathfrak{m} \otimes_k k'\) is a parabolic (resp. Borel) subalgebra of \(\mathfrak{a} \otimes_k k'\) .

This follows immediately from Prop. 14.

\section{SPLIT SEMI-SIMPLE LIE ALGEBRA DEFINED BY A REDUCED ROOT SYSTEM}

\subsection{FRAMED SEMI-SIMPLE LIE ALGEBRAS}

PROPOSITION 1. Let \((\mathfrak{g},\mathfrak{h})\) be a split semi-simple Lie algebra, R its root system, B a basis of R, and \((n(\alpha ,\beta))_{\alpha ,\beta \in \mathrm{B}}\) the corresponding Cartan matrix. For all \(\alpha \in \mathrm{B}\) , let \(X_{\alpha}\in \mathfrak{g}^{\alpha}\) , \(X_{-\alpha}\in \mathfrak{g}^{-\alpha}\) . Then, for \(\alpha ,\beta \in \mathrm{B}\) ,

\[
[ H _ {\alpha}, H _ {\beta} ] = 0\tag{1}
\]

\[
[ H _ {\alpha}, X _ {\beta} ] = n (\beta , \alpha) X _ {\beta}\tag{2}
\]

\[
[ H _ {\alpha}, X _ {- \beta} ] = - n (\beta , \alpha) X _ {- \beta}\tag{3}
\]

\[
[ X _ {- \alpha}, X _ {\beta} ] = 0 \quad \text { if } \alpha \neq \beta\tag{4}
\]

\[
(\operatorname{ad} X _ {\alpha}) ^ {1 - n (\beta , \alpha)} X _ {\beta} = 0 \quad \text { if } \alpha \neq \beta\tag{5}
\]

\[
(\operatorname{ad} X _ {- \alpha}) ^ {1 - n (\beta , \alpha)} X _ {- \beta} = 0 \quad \text { if } \alpha \neq \beta .\tag{6}
\]

The family \((H_{\alpha})_{\alpha \in \mathrm{B}}\) is a basis of \(\mathfrak{h}\) . If \(X_{\alpha} \neq 0\) and \(X_{-\alpha} \neq 0\) for all \(\alpha \in \mathrm{B}\) , the Lie algebra \(\mathfrak{g}\) is generated by the \(X_{\alpha}\) and the \(X_{-\alpha}\) ( \(\alpha \in \mathrm{B}\) ).

(Recall that, if \(\alpha, \beta \in \mathrm{B}\) and \(\alpha \neq \beta\) , \(n(\beta, \alpha)\) is an integer \(\leq 0\) , so formulas (5) and (6) make sense.)

Formulas (1), (2) and (3) are clear. If \(\alpha \neq \beta\) , \(\beta - \alpha\) is not a root since every element of \(R\) is a linear combination of elements of \(B\) with integer coefficients all of the same sign (Chap. VI, §1, no. 6, Th. 3). This proves (4). In view of Chap. VI, §1, no. 3, Prop. 9, this also proves that the \(\alpha\) -chain defined by \(\beta\) is

\[
\{\beta , \beta + \alpha , \dots , \beta - n (\beta , \alpha) \alpha \};
\]

hence \(\beta + (1 - n(\beta, \alpha))\alpha \notin \mathrm{R}\) , which proves (5). The equality (6) is established in a similar way. The family \((H_{\alpha})_{\alpha \in \mathrm{B}}\) is a basis of \(\mathrm{R}^{\vee}\) , and hence of \(\mathfrak{h}\) . If \(X_{\alpha} \neq 0\) and \(X_{-\alpha} \neq 0\) for all \(\alpha \in \mathrm{B}\) , then \([X_{\alpha}, X_{-\alpha}] = \lambda_{\alpha}H_{\alpha}\) with \(\lambda_{\alpha} \neq 0\) , so the last assertion follows from §3, no. 3, Prop. 9 (iii).

DEFINITION 1. Let \((\mathfrak{g},\mathfrak{h})\) be a split semi-simple Lie algebra, R its root system. A framing of \((\mathfrak{g},\mathfrak{h})\) is a pair \((\mathrm{B},(X_{\alpha})_{\alpha \in \mathrm{B}})\) , where B is a basis of R, and where, for all \(\alpha \in \mathrm{B}\) , \(X_{\alpha}\) is a non-zero element of \(\mathfrak{g}^{\alpha}\) . A framed semi-simple Lie algebra is a sequence \((\mathfrak{g},\mathfrak{h},\mathrm{B},(X_{\alpha})_{\alpha \in \mathrm{B}})\) where \((\mathfrak{g},\mathfrak{h})\) is a split semi-simple Lie algebra, and where \((\mathrm{B},(X_{\alpha})_{\alpha \in \mathrm{B}})\) is a framing of \((\mathfrak{g},\mathfrak{h})\) .

A framing of \(\mathfrak{g}\) is a framing of \((\mathfrak{g},\mathfrak{h})\) , where \(\mathfrak{h}\) is a splitting Cartan subalgebra of \(\mathfrak{g}\) .

Let \(a_1 = (\mathfrak{g}_1, \mathfrak{h}_1, \mathrm{B}_1, (X_\alpha^1)_{\alpha \in \mathrm{B}_1})\) and \(a_2 = (\mathfrak{g}_2, \mathfrak{h}_2, \mathrm{B}_2, (X_\alpha^2)_{\alpha \in \mathrm{B}_2})\) be framed semi-simple Lie algebras. An isomorphism from \(a_1\) to \(a_2\) is an isomorphism \(\varphi\) from \(\mathfrak{g}_1\) to \(\mathfrak{g}_2\) that takes \(\mathfrak{h}_1\) to \(\mathfrak{h}_2\) , \(\mathrm{B}_1\) to \(\mathrm{B}_2\) , and \(X_\alpha^1\) to \(X_{\psi \alpha}^2\) for all \(\alpha \in \mathrm{B}_1\) (where \(\psi\) is the contragredient map of \(\varphi | \mathfrak{h}_1\) ). In this case, \(\varphi\) is said to transform the framing \((\mathrm{B}_1, (X_\alpha^1)_{\alpha \in \mathrm{B}_1})\) to the framing \((\mathrm{B}_2, (X_\alpha^2)_{\alpha \in \mathrm{B}_2})\) .

If \((\mathrm{B},(X_{\alpha})_{\alpha\in\mathrm{B}})\) is a framing of \((\mathfrak{g},\mathfrak{h})\) , there exists, for all \(\alpha\in B\) , a unique element \(X_{-\alpha}\) of \(g^{-\alpha}\) such that \([X_{\alpha},X_{-\alpha}]=-H_{\alpha}\) (§2, no. 2, Th. 1 (iv)). The family \((X_{\alpha})_{\alpha\in\mathrm{B}\cup(-\mathrm{B})}\) is called the generating family defined by the framing (cf.~Prop.1). This is also the generating family defined by the framing \((-\mathrm{B},(X_{\alpha})_{\alpha\in-\mathrm{B}})\) . For all \(\alpha\in\mathrm{B}\cup(-\mathrm{B})\) , let \(t_{\alpha}\in k^{*}\) , and assume that \(t_{\alpha}t_{-\alpha}=1\) for all \(\alpha\in B\) . Then \((t_{\alpha}X_{\alpha})_{\alpha\in\mathrm{B}\cup(-\mathrm{B})}\) is the generating family defined by the framing \((\mathrm{B},(t_{\alpha}X_{\alpha})_{\alpha\in\mathrm{B}})\) .

\subsection{A PRELIMINARY CONSTRUCTION}

In this number and the next, we denote by R a reduced root system in a vector space V and by B a basis of R. We denote by \((n(\alpha,\beta))_{\alpha,\beta\in\mathbb{B}}\) the Cartan matrix relative to B. Recall that \(n(\alpha,\beta)=\langle\alpha,\beta^{\vee}\rangle\) . We are going to show that R is the root system of a split semi-simple Lie algebra which is unique up to isomorphism. In the main we shall be considering the Lie algebra defined by the relations in Prop. 1.

The construction in this number applies to any square matrix \((n(\alpha,\beta))_{\alpha,\beta\in\mathbb{B}}\) over k with non-zero determinant and such that \(n(\alpha,\alpha)=2\) for all \(\alpha\in B\) (cf.~Chap. VI, §1, no. 10, formula (14)).

Let E be the free associative algebra of the set B over k. Recall that E is N-graded (Algebra, Chap. III, §3, no. 1, Example 3). We are going to associate to each \(\alpha\in B\) endomorphisms \(X_{-\alpha}^{0}, H_{\alpha}^{0}, X_{\alpha}^{0}\) of the vector space E, of degrees 1, 0, -1 respectively. For any word \((\alpha_{1},\ldots,\alpha_{n})\) in elements of B, put

\[
X _ {- \alpha} ^ {0} (\alpha_ {1}, \ldots , \alpha_ {n}) = (\alpha , \alpha_ {1}, \ldots , \alpha_ {n})\tag{7}
\]

\[
H _ {\alpha} ^ {0} \left(\alpha_ {1}, \dots , \alpha_ {n}\right) = \left(- \sum_ {i = 1} ^ {n} n \left(\alpha_ {i}, \alpha\right)\right) \left(\alpha_ {1}, \dots , \alpha_ {n}\right).\tag{8}
\]

On the other hand, \(X_{\alpha}^{0}(\alpha_{1},\ldots ,\alpha_{n})\) is defined by induction on \(n\) using the formula

\[
X _ {\alpha} ^ {0} (\alpha_ {1}, \ldots , \alpha_ {n}) = (X _ {- \alpha_ {1}} ^ {0} X _ {\alpha} ^ {0} - \delta_ {\alpha , \alpha_ {1}} H _ {\alpha} ^ {0}) (\alpha_ {2}, \ldots , \alpha_ {n})\tag{9}
\]

where \(\delta_{\alpha, \alpha_1}\) is the Kronecker symbol; it is understood that \(X_{\alpha}^{0}(\alpha_{1}, \ldots, \alpha_{n})\) is zero if \((\alpha_{1}, \ldots, \alpha_{n})\) is the empty word.

Lemma 1. For all \(\alpha, \beta \in \mathbf{B}\) , we have

\[
[ X _ {\alpha} ^ {0}, X _ {- \alpha} ^ {0} ] = - H _ {\alpha} ^ {0}\tag{10}
\]

\[
[ H _ {\alpha} ^ {0}, H _ {\beta} ^ {0} ] = 0\tag{11}
\]

\[
[ H _ {\alpha} ^ {0}, X _ {\beta} ^ {0} ] = n (\beta , \alpha) X _ {\beta} ^ {0}\tag{12}
\]

\[
[ H _ {\alpha} ^ {0}, X _ {- \beta} ^ {0} ] = - n (\beta , \alpha) X _ {- \beta} ^ {0}\tag{13}
\]

\[
[ X _ {\alpha} ^ {0}, X _ {- \beta} ^ {0} ] = 0 \text {   if   } \alpha \neq \beta .\tag{14}
\]

Indeed, relation (9) can be written

\[
(X _ {\alpha} ^ {0} X _ {- \alpha_ {1}} ^ {0}) (\alpha_ {2}, \ldots , \alpha_ {n}) = (X _ {- \alpha_ {1}} ^ {0} X _ {\alpha} ^ {0}) (\alpha_ {2}, \ldots , \alpha_ {n}) - \delta_ {\alpha , \alpha_ {1}} H _ {\alpha} ^ {0} (\alpha_ {2}, \ldots , \alpha_ {n})
\]

which proves (10) and (14). Relation (11) is clear. Next

\[
\begin{array}{r l} [ H _ {\alpha} ^ {0}, X _ {- \beta} ^ {0} ] (\alpha_ {1}, \ldots , \alpha_ {n}) & = H _ {\alpha} ^ {0} (\beta , \alpha_ {1}, \ldots , \alpha_ {n}) + \left(\sum_ {i = 1} ^ {n} n (\alpha_ {i}, \alpha)\right) (\beta , \alpha_ {1}, \ldots , \alpha_ {n}) \\ & = - n (\beta , \alpha) (\beta , \alpha_ {1}, \ldots , \alpha_ {n}) \\ & = - n (\beta , \alpha) X _ {- \beta} ^ {0} (\alpha_ {1}, \ldots , \alpha_ {n}) \end{array}
\]

hence (13). Finally,

\[
\begin{array}{l} 0 = [ H _ {\alpha} ^ {0}, [ X _ {\beta} ^ {0}, X _ {- \gamma} ^ {0} ] ] \quad \text { by (10), (11), (14)} \\ \qquad = [ [ H _ {\alpha} ^ {0}, X _ {\beta} ^ {0} ], X _ {- \gamma} ^ {0} ] + [ X _ {\beta} ^ {0}, [ H _ {\alpha} ^ {0}, X _ {- \gamma} ^ {0} ] ] \\ \qquad = [ [ H _ {\alpha} ^ {0}, X _ {\beta} ^ {0} ] - n (\gamma , \alpha) X _ {\beta} ^ {0}, X _ {- \gamma} ^ {0} ] \quad \text { by (13)} \\ \qquad = [ [ H _ {\alpha} ^ {0}, X _ {\beta} ^ {0} ] - n (\beta , \alpha) X _ {\beta} ^ {0}, X _ {- \gamma} ^ {0} ] \quad \text { by (14); } \end{array}\tag{15}
\]

now, considering the empty word immediately gives

\[
([ H _ {\alpha} ^ {0}, X _ {\beta} ^ {0} ] - n (\beta , \alpha) X _ {\beta} ^ {0}) (\emptyset) = 0
\]

so (15) implies that

\[
\left(\left[ H _ {\alpha} ^ {0}, X _ {\beta} ^ {0} \right] - n (\beta , \alpha) X _ {\beta} ^ {0}\right) X _ {- \gamma_ {1}} ^ {0} X _ {- \gamma_ {2}} ^ {0} \dots X _ {- \gamma_ {n}} ^ {0} (\varnothing) = 0
\]

for all \(\gamma_1, \ldots, \gamma_n \in \mathrm{B}\) ; this proves (12).

Lemma 2. The endomorphisms \(X_{\alpha}^{0}\) , \(H_{\beta}^{0}\) , \(X_{-\gamma}^{0}\) , where \(\alpha, \beta, \gamma \in \mathrm{B}\) , are linearly independent.

Since \(X_{-\alpha}^{0}(\varnothing) = \alpha\) , it is clear that the \(X_{-\alpha}^{0}\) are linearly independent. Assume that \(\sum_{\alpha} a_{\alpha} H_{\alpha}^{0} = 0\) ; then, for all \(\beta \in B\) ,

\[
0 = \left[ \sum_ {\alpha} a _ {\alpha} H _ {\alpha} ^ {0}, X _ {- \beta} ^ {0} \right] = - \sum_ {\alpha} a _ {\alpha} n (\beta , \alpha) X _ {- \beta} ^ {0};
\]

since \(\operatorname{det}(n(\beta, \alpha)) \neq 0\) , it follows that \(a_{\alpha} = 0\) for all \(\alpha\) . Assume that \(\sum_{\alpha} a_{\alpha} X_{\alpha}^{0} = 0\) . In view of formulas (7), (8), (9),

\[
X _ {\alpha} ^ {0} (\beta) = 0,
\]

\[
X _ {\alpha} ^ {0} (\beta , \beta) = 2 \delta_ {\alpha \beta} \beta
\]

for all \(\beta \in \mathrm{B}\) . It follows that \(a_{\beta} = 0\) for all \(\beta\) . Since \(X_{\alpha}^{0}\) , \(H_{\alpha}^{0}\) , \(X_{-\alpha}^{0}\) are of degree \(-1, 0, 1\) , respectively, the lemma follows from what has gone before.

Let I be the set \(B \times \{-1,0,1\}\) . Put \(x_{\alpha} = (\alpha, -1)\) , \(h_{\alpha} = (\alpha, 0)\) , and \(x_{-\alpha} = (\alpha, 1)\) . Let a be the Lie algebra defined by the generating family I and the following set R of relators:

\[
\begin{array}{l} \left[ h _ {\alpha}, h _ {\beta} \right] \\ \left[ h _ {\alpha}, x _ {\beta} \right] - n (\beta , \alpha) x _ {\beta} \\ \left[ h _ {\alpha}, x _ {- \beta} \right] + n (\beta , \alpha) x _ {- \beta} \\ \left[ x _ {\alpha}, x _ {- \alpha} \right] + h _ {\alpha} \\ \left[ x _ {\alpha}, x _ {- \beta} \right] \quad \text {if} \alpha \neq \beta \end{array}
\]

(cf.~Chap. II, §2, no. 3). By Lemma 1, there exists a unique linear representation \(\rho\) of \(\mathfrak{a}\) on \(\mathrm{E}\) such that

\[
\rho (x _ {\alpha}) = X _ {\alpha} ^ {0}, \rho (h _ {\alpha}) = H _ {\alpha} ^ {0}, \rho (x _ {- \alpha}) = X _ {- \alpha} ^ {0}.
\]

In view of Lemma 2, this proves the following result:

Lemma 3. The canonical images in \(\mathfrak{a}\) of the elements \(x_{\alpha}, h_{\beta}, x_{-\gamma}\) , where \(\alpha, \beta, \gamma \in \mathrm{B}\) , are linearly independent.

In the following, we identify \(x_{\alpha}\) , \(h_{\alpha}\) , \(x_{-\alpha}\) with their canonical images in a.

Lemma 4. There exists a unique involutive automorphism \(\theta\) of \(\mathfrak{a}\) such that

\[
\theta (x _ {\alpha}) = x _ {- \alpha}, \theta (x _ {- \alpha}) = x _ {\alpha}, \theta (h _ {\alpha}) = - h _ {\alpha}
\]

for all \(\alpha\in B\) .

Indeed, there exists an involutive automorphism of the free Lie algebra \(L(I)\) satisfying these conditions. It leaves \(R \cup (-R)\) stable, and hence defines by passage to the quotient an involutive automorphism of a satisfying the conditions of the lemma. The uniqueness follows from the fact that a is generated by the elements \(x_{\alpha}\) , \(h_{\alpha}\) , \(x_{-\alpha}\) ( \(\alpha \in B\) ).

This automorphism is called the canonical involutive automorphism of a.

Let Q be the set of radical weights of R; this is a free Z-module with basis B (Chap. VI, §1, no. 9). There exists a graduation of type Q on the free Lie algebra L(I) such that \(x_{\alpha}\) , \(h_{\alpha}\) , \(x_{-\alpha}\) are of degrees \(\alpha\) , 0, - \(\alpha\) , respectively (Chap. II, §2, no. 6). Now the elements of R are homogeneous. Hence there exists a unique graduation of type Q on a compatible with the Lie algebra structure of a and such that \(x_{\alpha}\) , \(h_{\alpha}\) , \(x_{-\alpha}\) are of degrees \(\alpha\) , 0, - \(\alpha\) , respectively. For any \(\mu \in Q\) , denote by \(a^{\mu}\) the set of elements of a homogeneous of degree \(\mu\) .

Lemma 5. Let \(z \in \mathfrak{a}\) . Then \(z \in \mathfrak{a}^{\mu}\) if and only if \([h_{\alpha}, z] = \langle \mu, \alpha^{\vee} \rangle z\) for all \(\alpha \in \mathrm{B}\) .

For \(\mu \in \mathrm{Q}\) , let \(\mathfrak{a}^{(\mu)}\) be the set of \(x \in \mathfrak{a}\) such that \([h_{\alpha}, x] = \langle \mu, \alpha^{\vee} \rangle x\) for all \(\alpha \in \mathrm{B}\) . The sum of the \(\mathfrak{a}^{(\mu)}\) is direct. To prove the lemma, it therefore suffices to show that \(\mathfrak{a}^{\mu} \subset \mathfrak{a}^{(\mu)}\) . Let \(\alpha \in \mathrm{B}\) . The endomorphism \(u\) of the vector space \(\mathfrak{a}\) such that \(u|\mathfrak{a}^{\mu} = \langle \mu, \alpha^{\vee} \rangle\) .1 is a derivation of \(\mathfrak{a}\) such that \(ux = (\operatorname{ad} h_{\alpha})\) . \(x\) for \(x = x_{\beta}\) , \(x = h_{\beta}\) , \(x = x_{-\beta}\) ; hence \(u = \operatorname{ad} h_{\alpha}\) , which proves our assertion.

Remark. It follows from Lemma 5 that every ideal of \(\mathfrak{a}\) is homogeneous, since it is stable under the ad \(h_{\alpha}\) .

Denote by \(Q_{+}\) (resp. \(Q_{-}\) ) the set of linear combinations of elements of B with positive (resp. negative) integer coefficients, not all zero. Put \(a_{+} = \sum_{\mu \in Q_{+}} a^{\mu}\) and \(a_{-} = \sum_{\mu \in Q_{-}} a^{\mu}\) . Since \(Q_{+} + Q_{+} \subset Q_{+}\) and \(Q_{-} + Q_{-} \subset Q_{-}\) , \(a_{+}\) and \(a_{-}\) are Lie subalgebras of a.

PROPOSITION 2. (i) The Lie algebra \(\mathfrak{a}_{+}\) is generated by the family \((x_{\alpha})_{\alpha \in \mathbb{B}}\) .

\begin{enumerate}
\def\labelenumi{(\roman{enumi})}
\setcounter{enumi}{1}
\item
  The Lie algebra \(\mathfrak{a}_{-}\) is generated by the family \((x_{-\alpha})_{\alpha \in \mathrm{B}}\) .
\item
  The family \((h_{\alpha})_{\alpha \in \mathrm{B}}\) is a basis of the vector space \(\mathfrak{a}^0\) .
\item
  The vector space \(\mathfrak{a}\) is the direct sum of \(\mathfrak{a}_{+},\mathfrak{a}^{0},\mathfrak{a}_{-}\) .
\end{enumerate}

Let \(\mathfrak{r}\) (resp. \(\mathfrak{n}\) ) be the Lie subalgebra of \(\mathfrak{a}\) generated by \((x_{\alpha})_{\alpha \in \mathrm{B}}\) (resp. \((x_{-\alpha})_{\alpha \in \mathrm{B}}\) ), and \(\mathfrak{h}\) the vector subspace of \(\mathfrak{a}\) generated by \((h_{\alpha})_{\alpha \in \mathrm{B}}\) . Since the \(x_{\alpha}\) are homogeneous elements of \(\mathfrak{a}_{+}\) , \(\mathfrak{r}\) is a graded subalgebra of \(\mathfrak{a}_{+}\) ; hence, \([\mathfrak{h},\mathfrak{r}]\subset \mathfrak{r}\) , so \(\mathfrak{h} + \mathfrak{r}\) is a subalgebra of \(\mathfrak{a}\) ; since

\[
[ x _ {- \alpha}, x _ {\beta} ] = \delta_ {\alpha \beta} h _ {\alpha},
\]

\([x_{-\alpha}, \mathfrak{r}] \subset \mathfrak{h} + \mathfrak{r}\) for all \(\alpha \in B\) . Similarly, \(\mathfrak{n}\) is a graded subalgebra of \(\mathfrak{a}_{-}\) , one has \([\mathfrak{h}, \mathfrak{n}] \subset \mathfrak{n}\) , \(\mathfrak{h} + \mathfrak{n}\) is a subalgebra of \(\mathfrak{n}\) , and \([x_{\alpha}, \mathfrak{n}] \subset \mathfrak{h} + \mathfrak{n}\) for all \(\alpha \in B\) . Put \(a' = r + h + n\) . The preceding shows that \(a'\) is stable under ad \(x_{\alpha}\) , ad \(h_{\alpha}\) and ad \(x_{-\alpha}\) for all \(\alpha \in B\) , and hence is an ideal of a. Since \(a'\) contains \(x_{\alpha}\) , \(h_{\alpha}\) , \(x_{-\alpha}\) for all \(\alpha \in B\) , \(a' = a\) . It follows from this that the inclusions \(r \subset a_{+}\) , \(h \subset a^{0}\) , \(n \subset a_{-}\) are equalities, which proves the proposition.

PROPOSITION 3. The Lie algebra \(\mathfrak{a}_{+}\) (resp. \(\mathfrak{a}_{-}\) ) is a free Lie algebra with basic family \((x_{\alpha})_{\alpha \in \mathrm{B}}\) (resp. \((x_{-\alpha})_{\alpha \in \mathrm{B}}\) ) (cf.~Chap. II, §2, no. 3).

Let L be the Lie subalgebra of E generated by B. By Chap. II, §3, Th. 1, L can be identified with the free Lie algebra generated by B. The left regular representation of E on itself is clearly injective, and defines by restriction to L an injective representation \(\rho'\) of the Lie algebra L on E. Let \(\varphi\) be the unique homomorphism from L to \(\mathfrak{a}_{-}\) which takes \(\alpha\) to \(x_{-\alpha}\) for all \(\alpha \in B\) . Then, for all \(\alpha \in B\) , \(\rho(\varphi(\alpha))\) is the endomorphism of left multiplication by \(\alpha\) on E, so \(\rho \circ \varphi = \rho'\) , which proves that \(\varphi\) is injective. Thus, \((x_{-\alpha})_{\alpha \in B}\) is a basic family for \(\mathfrak{a}_{-}\) . Since \(\theta(x_{-\alpha}) = x_{\alpha}\) for all \(\alpha\) (cf.~Lemma 4), \((x_{\alpha})_{\alpha \in B}\) is a basic family for \(\mathfrak{a}_{+}\) .